\chapter{Metodologia e \textit{Framework} Proposto}
\label{cap:metodologia}

Este capítulo está organizado em duas partes. A primeira, apresentada na Seção~\ref{sec:validacao-empirica}, descreve a avaliação empírica e a evolução do protocolo experimental. O estudo compreende uma fase preliminar, executada por meio de APIs externas, e uma fase local consolidada. A fase local foi conduzida em duas etapas sequenciais. Na etapa inicial, cinco estratégias não invasivas de tratamento de contexto foram comparadas em cinco modelos e dois tamanhos de contexto, formando a matriz de referência para a escolha da estratégia base. Esses resultados foram concebidos para complementar a avaliação via API, cuja execução foi limitada por restrições na implementação e na customização da janela de contexto descritas na Seção~\ref{subsec:evolucao-protocolo}. A etapa de ampliação foi motivada pela hipótese de que a seleção de trechos de contexto pode ser formulada como um problema de otimização, formalizado na Seção~\ref{sec:framework}. Nessa etapa, LLMLingua-GPT2, LLMLingua-2, Selective Context e CPC-MiniLM foram incorporados como quatro técnicas adicionais, ampliando o conjunto comparativo para nove técnicas nas dez combinações de modelo e orçamento. A segunda parte do capítulo, apresentada na Seção~\ref{sec:framework}, formaliza a arquitetura do \textit{framework} proposto, que descreve a compressão seletiva de contexto como um problema da mochila de múltipla escolha.

\section{Validação Empírica das Estratégias Não Invasivas}
\label{sec:validacao-empirica}

\subsection{Objetivo e Desenho Experimental}
\label{subsec:desenho}

A validação empírica tem por objetivo medir quanto de informação cada técnica de tratamento de contexto preserva, isto é, em que medida o modelo ainda recupera a resposta correta após a redução ou a segmentação da entrada, em comparação com o \textit{baseline} sem tratamento. A questão central não é qual modelo apresenta melhor desempenho absoluto e sim como o compromisso entre qualidade e tempo de execução varia entre técnicas, famílias de modelos e tamanhos de contexto. Para responder a essa questão sem alterar retrospectivamente o protocolo já executado, a matriz local foi construída em duas etapas sequenciais, descritas a seguir.

\paragraph{Etapa inicial.}
Na etapa inicial, cinco estratégias não invasivas de tratamento de contexto foram comparadas, formando a matriz local inicial e fornecendo a referência para a escolha da estratégia base. As cinco estratégias são descritas a seguir.

\begin{itemize}[leftmargin=1.6em]
\item \textbf{\textit{Raw}}. Estratégia de referência (\textit{baseline}), em que o contexto completo é enviado ao modelo sem qualquer tratamento.

\item \textbf{\textit{Parallel Window}}. Segmentação do contexto em janelas processadas em paralelo, com agregação posterior das respostas parciais \cite{ratner2023parallelwindows}.

\item \textbf{\textit{Sliding Window}}. Segmentação com janelas deslizantes e sobreposição, preservando continuidade semântica entre segmentos sucessivos.

\item \textbf{RIG}. Seleção de trechos por ganho de informação relevante, com ordenação do tipo \textit{Dartboard}, priorizando o conteúdo que mais reduz a incerteza sobre a resposta \cite{pickett2024rig}.

\item \textbf{\textit{Semantic Compression}}. Compressão semântica orientada por instruções de \textit{prompt}, sem modelo de compressão adicional \cite{gilbert2023semanticcompression}.
\end{itemize}

Essa etapa foi executada integralmente nos cinco modelos e nos dois orçamentos apresentados a seguir.

\paragraph{Etapa de ampliação.}
Após a conclusão da etapa inicial, quatro compressores foram acrescentados como técnicas adicionais, preservando os mesmos modelos, orçamentos, \textit{benchmarks} e critérios de agregação e permitindo apresentar as nove técnicas nas mesmas tabelas.

\begin{itemize}[leftmargin=1.6em]
\item \textbf{LLMLingua-GPT2}. \textit{Baseline} externo de compressão, utilizado como referência para avaliar o ganho obtido pelos compressores mais recentes.

\item \textbf{LLMLingua-2}. Candidato ao conjunto de opções do MCKP, avaliado como técnica de compressão baseada em modelo.

\item \textbf{\textit{Selective Context}}. Candidato ao conjunto de opções do MCKP, avaliado como técnica de seleção de trechos por relevância.

\item \textbf{CPC-MiniLM}. Candidato ao conjunto de opções do MCKP, avaliado como técnica de compressão orientada por similaridade semântica.
\end{itemize}

Os quatro compressores foram executados nas dez combinações de modelo e orçamento. Uma chamada do CPC-MiniLM ao DeepSeek-R1 32B, no orçamento de 8192 \textit{tokens}, terminou por \textit{timeout}; essa observação é excluída das agregações dependentes da resposta e identificada nas tabelas. As implementações e adaptações dos quatro compressores são detalhadas na Seção~\ref{subsec:implementacao}.


\paragraph{Modelos avaliados.}
Os modelos foram selecionados para cobrir eixos distintos de projeto, e não por serem os de maior porte. A Tabela~\ref{tab:modelos} apresenta os cinco modelos e seu papel na matriz de avaliação, contemplando arquitetura densa em escala pequena e grande, arquitetura de mistura de especialistas esparsa e modelos ajustados para raciocínio. A escolha por famílias distintas, que reúnem Meta, Google, Alibaba, DeepSeek e OpenAI, busca reduzir a dependência do resultado em relação a um único estilo de treinamento.

\begin{table}[htbp]
\centering
\caption{Modelos avaliados e diversidade arquitetural da matriz de \textit{benchmarks}}
\label{tab:modelos}
\small
\setlength{\tabcolsep}{5pt}
\begin{tabular}{p{3.3cm}p{2.6cm}p{2.6cm}p{3.9cm}}
\toprule
\textbf{Modelo} & \textbf{Família} & \textbf{Arquitetura} & \textbf{Papel na matriz} \\
\midrule
\texttt{llama3.1:8b} & Llama 3.1 (Meta) & Densa, 8B & \textit{Baseline} pequeno e rápido \\
\addlinespace
\texttt{gemma4:26b} & Gemma 4 (Google) & Densa, 26B & Densa de grande porte, família distinta de Llama \\
\addlinespace
\texttt{qwen3:30b-a3b} & Qwen 3 (Alibaba) & MoE esparsa, 30B com 3B ativos & Mistura de especialistas, ativa apenas parte dos pesos \\
\addlinespace
\texttt{deepseek-r1:32b} & DeepSeek-R1 & Raciocínio, 32B & Ajustado para cadeia de raciocínio \\
\addlinespace
\texttt{gpt-oss:20b} & GPT-OSS (OpenAI) & MoE, 20B & Família aberta da OpenAI, estilo de treinamento distinto \\
\bottomrule
\end{tabular}
\end{table}

\paragraph{Orçamentos de contexto.}
A avaliação de cada modelo considera dois orçamentos de contexto, de 8192 e de 32768 \textit{tokens}, correspondentes aos extremos de pressão de compressão. O orçamento de 8192 \textit{tokens} impõe forte truncamento, exigindo que a estratégia selecione bem o conteúdo mantido, enquanto o orçamento de 32768 \textit{tokens} mede o teto de qualidade, uma vez que quase todo o conteúdo é acomodado. Uma calibração preliminar mais ampla, conduzida com o modelo \textit{baseline}, indicou que os contextos de 16384 e de 32768 \textit{tokens} saturam em um mesmo patamar de acurácia, próximo de 0.954, de modo que o ponto intermediário não acrescenta informação relevante à curva de compromisso e foi omitido para os demais modelos. O modelo \textit{baseline} preserva os três contextos como curva completa de referência.

\paragraph{Suíte de \textit{benchmarks}.}
A suíte de \textit{benchmarks} reúne sete tarefas de contexto longo restritas a texto, agrupadas em quatro categorias. O LongBench \cite{bai2023longbench} e o ZeroSCROLLS \cite{shaham2023zeroscrolls} reúnem múltiplas tarefas de contexto longo em formato de avaliação sem ajuste específico, cobrindo compreensão, sumarização e resposta a perguntas. A resposta a perguntas de múltiplos saltos é avaliada pelo HotpotQA \cite{yang2018hotpotqa} e pelo MuSiQue \cite{trivedi2022musique}, que exigem combinar evidências dispersas em vários documentos e, por isso, penalizam compressões que descartam fatos de ligação. A resposta a perguntas de domínio aberto é avaliada pelo Natural Questions \cite{kwiatkowski2019natural} e pelo TriviaQA \cite{joshi2017triviaqa}, que utilizam consultas reais de usuários e questões sobre documentos. A sumarização orientada por consulta é avaliada pelo QMSum \cite{zhong2021qmsum}, que requer selecionar e resumir trechos relevantes ao longo de diálogos extensos. \textit{Benchmarks} puramente sintéticos, como o Needle-in-a-Haystack e o RULER, foram excluídos da matriz principal pelas razões discutidas na Seção~\ref{subsec:revisao-benchmarks}. Para cada combinação de modelo, orçamento de contexto, estratégia e \textit{benchmark}, registram-se a acurácia, a latência e a taxa de compressão efetiva. Valores não disponíveis são indicados por \textit{n/d}.

\paragraph{Artefatos e reprodutibilidade.}
As decisões metodológicas relativas à seleção dos modelos, aos orçamentos de contexto e à composição da suíte estão registradas em \texttt{tests/README.md} e em \texttt{tests/ollama-local/README\_DECISAO\_TESTES.md}. Nos dois orçamentos comparáveis, a etapa inicial da matriz local produziu 1100 resultados individuais, correspondentes a 22 casos para cada uma das 50 combinações de modelo, orçamento e estratégia. Esses resultados correspondem a 350 células agregadas por \textit{benchmark}. Os artefatos de execução incluem os \textit{logs}, os resultados por caso nos formatos CSV e JSON e as tabelas agregadas disponíveis em \texttt{tests/ollama-local/benchmark\_matrix\_artigo}. A execução adicional do Llama 3.1 8B com 16384 \textit{tokens} foi preservada como ponto intermediário de calibração, mas não integra as tabelas comparativas dos dois orçamentos extremos.

\subsection{Evolução do Protocolo Experimental}
\label{subsec:evolucao-protocolo}

O desenho experimental consolidado na Seção~\ref{subsec:desenho} resulta de uma evolução do protocolo de avaliação conduzida em duas fases. A fase preliminar utilizou modelos servidos por provedores de servidores externos, especificamente a Groq e a Cerebras, o que permitiu avaliar modelos de grande porte, como o Llama 3.3 70B e o GPT-OSS 120B, em um conjunto inicial de \textit{benchmarks} composto por Needle-in-a-Haystack, RULER e LongBench. Essa configuração aproveitou o acesso gratuito então oferecido por esses provedores.

Ao longo do trabalho, contudo, as condições de acesso foram alteradas. Os planos gratuitos da Groq e da Cerebras foram revisados, diversos modelos foram descontinuados ou indisponibilizados, e os limites de requisição passaram a restringir a execução da matriz completa. Essa alteração faz parte do processo de pesquisa apoiado em infraestrutura de terceiros, cuja disponibilidade não é controlada pelo autor. Para preservar a reprodutibilidade e a autonomia da avaliação, a segunda fase migrou para modelos executados localmente por meio do Ollama, o que motivou a matriz de cinco modelos de arquiteturas distintas nos orçamentos de 8192 e de 32768 \textit{tokens} descrita na Seção~\ref{subsec:desenho}. Dentro dessa fase local, as cinco estratégias originais foram avaliadas primeiro; os quatro compressores adicionais foram incorporados posteriormente, sem alterar os modelos, os orçamentos, a suíte ou os critérios de agregação.

\subsection{Fase Preliminar, Avaliação via APIs Externas}
\label{subsec:fase-preliminar}

Os resultados da fase preliminar são reportados a seguir e preservados no texto por transparência do processo de pesquisa. Essa fase avaliou quatro modelos, o Llama 3.1 8B, o Llama 3.3 70B, o GPT-OSS 20B e o GPT-OSS 120B, nos \textit{benchmarks} Needle-in-a-Haystack, RULER e LongBench, registrando a acurácia e a latência média de resposta por combinação de modelo, estratégia e \textit{benchmark}. A coluna \textbf{Média} corresponde à média das métricas nos \textit{benchmarks} avaliados. Valores não disponíveis, decorrentes sobretudo dos limites de requisição descritos, são indicados por \textit{n/d}.

\subsection{Resultados de Acurácia da Fase Preliminar}
\label{subsec:resultados-acuracia}

A Tabela~\ref{tab:acuracia} apresenta a acurácia obtida por cada estratégia nos três \textit{benchmarks}, agrupada por modelo. Em todas as tabelas de resultados deste capítulo, valores em \textbf{negrito} destacam o melhor resultado por coluna dentro de cada modelo e valores \underline{sublinhados} o segundo melhor; nas tabelas de latência, menor é melhor. A marcação \textit{n/d} indica ausência de dados na configuração correspondente.

\begin{table}[htbp]
\centering
\caption{Acurácia por estratégia, modelo e \textit{benchmark} de contexto longo}
\label{tab:acuracia}
\small
\setlength{\tabcolsep}{4.5pt}
\begin{tabular}{lccc|c}
\toprule
& \multicolumn{3}{c|}{\textbf{Benchmarks de Contexto Longo}} & \\
\cmidrule(lr){2-4}
\textbf{Estratégia} & \textbf{Needle} & \textbf{RULER} & \textbf{LongBench} & \textbf{Média} \\
\midrule
\multicolumn{5}{l}{\textit{Llama 3.1 8B}} \\
\textit{Parallel Window}      & \textbf{1.00}    & \textbf{1.00}    & 0.94         & \textbf{0.98} \\
\textit{Raw}                  & \textbf{1.00}    & \underline{0.94} & 0.94         & \underline{0.96} \\
RIG                           & \textbf{1.00}    & \textbf{1.00}    & 0.94         & \textbf{0.98} \\
\textit{Semantic Compression} & \underline{0.33} & 0.56             & 0.94         & 0.61 \\
\textit{Sliding Window}       & \textbf{1.00}    & \textbf{1.00}    & 0.94         & \textbf{0.98} \\
\midrule
\multicolumn{5}{l}{\textit{Llama 3.3 70B}} \\
\textit{Parallel Window}      & \textbf{1.00}    & \textbf{1.00}    & \underline{0.94} & \underline{0.98} \\
\textit{Raw}                  & \textbf{1.00}    & \textbf{1.00}    & \textbf{1.00}    & \textbf{1.00} \\
RIG                           & \textbf{1.00}    & \textbf{1.00}    & \underline{0.94} & \underline{0.98} \\
\textit{Semantic Compression} & \underline{0.33} & \underline{0.69} & 0.38             & 0.47 \\
\textit{Sliding Window}       & \textbf{1.00}    & \textbf{1.00}    & \underline{0.94} & \underline{0.98} \\
\midrule
\multicolumn{5}{l}{\textit{GPT-OSS 20B}} \\
\textit{Parallel Window}      & \textbf{1.00}    & 0.69             & \textit{n/d} & 0.85 \\
\textit{Raw}                  & \textbf{1.00}    & \textbf{0.78}    & \textbf{0.50} & 0.76 \\
RIG                           & \textbf{1.00}    & 0.50             & \textit{n/d} & 0.75 \\
\textit{Semantic Compression} & \underline{0.73} & 0.00             & \textit{n/d} & 0.37 \\
\textit{Sliding Window}       & \textbf{1.00}    & \underline{0.72} & \textit{n/d} & \textbf{0.86} \\
\midrule
\multicolumn{5}{l}{\textit{GPT-OSS 120B}} \\
\textit{Parallel Window}      & \textbf{1.00}    & \underline{0.58} & \underline{0.00} & 0.53 \\
\textit{Raw}                  & \textbf{1.00}    & 0.56             & \textbf{0.17}    & \underline{0.58} \\
RIG                           & \underline{0.97} & 0.44             & \textit{n/d}     & \textbf{0.71} \\
\textit{Semantic Compression} & \textbf{1.00}    & 0.06             & \textit{n/d}     & 0.53 \\
\textit{Sliding Window}       & 0.87             & \textbf{0.64}    & \underline{0.00} & 0.50 \\
\bottomrule
\end{tabular}
\end{table}

Os resultados de acurácia indicam que as estratégias de segmentação, tanto \textit{Parallel Window} quanto \textit{Sliding Window}, e a seleção por RIG mantêm desempenho elevado nos modelos Llama, com média próxima de 0.98. A estratégia \textit{Raw} apresenta o melhor resultado isolado no Llama 3.3 70B, com média de 1.00, o que sugere que modelos maiores conseguem aproveitar o contexto completo de forma mais eficaz quando não há restrição de orçamento. A \textit{Semantic Compression} apresenta o pior desempenho de acurácia entre as estratégias avaliadas, com quedas acentuadas no \textit{benchmark} RULER, o que é compatível com a perda de informação associada a uma compressão agressiva e uniforme. Nos modelos GPT-OSS, observa-se degradação geral de acurácia no RULER e no LongBench, com maior dispersão entre as estratégias.

\subsection{Resultados de Latência da Fase Preliminar}
\label{subsec:resultados-latencia}

A Tabela~\ref{tab:latencia} apresenta a latência média de resposta, em milissegundos, por estratégia, modelo e \textit{benchmark}.

\begin{table}[htbp]
\centering
\caption{Latência média de resposta (ms) por estratégia, modelo e \textit{benchmark}}
\label{tab:latencia}
\small
\setlength{\tabcolsep}{4.5pt}
\begin{tabular}{lccc|c}
\toprule
& \multicolumn{3}{c|}{\textbf{Benchmarks de Contexto Longo}} & \\
\cmidrule(lr){2-4}
\textbf{Estratégia} & \textbf{Needle} & \textbf{RULER} & \textbf{LongBench} & \textbf{Média} \\
\midrule
\multicolumn{5}{l}{\textit{Llama 3.1 8B}} \\
\textit{Parallel Window}      & 6200          & \underline{3700} & 2400          & 4100  \\
\textit{Raw}                  & \textbf{318}  & \textbf{1000}    & \underline{344} & \textbf{554} \\
RIG                           & 6300          & 3900             & 6300          & 5500  \\
\textit{Semantic Compression} & 11700         & 7300             & 11700         & 10200 \\
\textit{Sliding Window}       & \underline{2300} & 3800          & \textbf{286}  & \underline{2129} \\
\midrule
\multicolumn{5}{l}{\textit{Llama 3.3 70B}} \\
\textit{Parallel Window}      & 271           & 2300          & 346          & 972  \\
\textit{Raw}                  & \textbf{247}  & \textbf{280}  & \textbf{316} & \textbf{281} \\
RIG                           & 3200          & 2500          & 1500         & 2400 \\
\textit{Semantic Compression} & 3600          & 5100          & 965          & 3222 \\
\textit{Sliding Window}       & \underline{262} & \underline{517} & \underline{336} & \underline{372} \\
\midrule
\multicolumn{5}{l}{\textit{GPT-OSS 20B}} \\
\textit{Parallel Window}      & 7134          & 3489             & \textit{n/d}  & 5312 \\
\textit{Raw}                  & \textbf{2775} & \underline{298}  & \textbf{2766} & \textbf{1946} \\
RIG                           & \underline{7055} & 2766          & \textit{n/d}  & \underline{4911} \\
\textit{Semantic Compression} & 14676         & \textbf{130}     & \textit{n/d}  & 7403 \\
\textit{Sliding Window}       & 7527          & 3399             & \textit{n/d}  & 5463 \\
\midrule
\multicolumn{5}{l}{\textit{GPT-OSS 120B}} \\
\textit{Parallel Window}      & \underline{7165} & 3372          & 159          & \underline{3565} \\
\textit{Raw}                  & \textbf{3507} & \underline{2141} & 231          & \textbf{1960} \\
RIG                           & 7272          & 3085             & \textit{n/d} & 5179 \\
\textit{Semantic Compression} & 16670         & \textbf{2071}    & \textit{n/d} & 9371 \\
\textit{Sliding Window}       & 7681          & 3638             & \textbf{71}  & 3797 \\
\bottomrule
\end{tabular}
\end{table}

Os resultados de latência mostram que a estratégia \textit{Raw} é consistentemente a mais rápida, o que é esperado, pois não há custo adicional de segmentação, seleção ou compressão. Entre as estratégias que efetivamente tratam o contexto, a \textit{Sliding Window} apresenta a menor latência no Llama 3.1 8B, com média de aproximadamente 2,1 segundos, enquanto RIG e \textit{Semantic Compression} apresentam as maiores latências, associadas ao custo de seleção por ganho de informação e ao custo de compressão semântica. Nos modelos GPT-OSS, as latências são substancialmente maiores para todas as estratégias baseadas em segmentação e compressão, com a \textit{Semantic Compression} atingindo os maiores tempos observados.

\subsection{Discussão da Fase Preliminar}
\label{subsec:discussao-fase-preliminar}

A análise conjunta das Tabelas~\ref{tab:acuracia} e \ref{tab:latencia} evidencia um compromisso entre acurácia e latência. A estratégia \textit{Raw} oferece a menor latência, porém sem qualquer mecanismo de tratamento de contexto, o que limita sua aplicabilidade em cenários que excedem a janela nominal do modelo. A \textit{Semantic Compression}, por outro lado, combina baixa acurácia com alta latência nos cenários avaliados, o que a torna pouco atrativa como estratégia isolada e uniforme. Entre as estratégias que tratam o contexto, a \textit{Sliding Window} apresenta o melhor equilíbrio na fase preliminar, mantendo acurácia média de 0.98 no Llama 3.1 8B, equivalente à das demais estratégias de segmentação e seleção, com a menor latência entre elas.

Um ponto que atravessa os dois conjuntos de resultados é que nenhuma estratégia uniforme domina todas as demais em todos os modelos e \textit{benchmarks}. Esse comportamento é consistente com a hipótese central desta dissertação, segundo a qual o tratamento uniforme do \textit{prompt} desperdiça orçamento em regiões redundantes e compromete regiões de alto valor informacional. A fase preliminar também expôs uma limitação do conjunto de \textit{benchmarks} então utilizado, discutida a seguir, que motivou a revisão da suíte de avaliação adotada na fase local.

\subsection{Revisão da Suíte de Benchmarks}
\label{subsec:revisao-benchmarks}

A fase preliminar apoiou-se em dois \textit{benchmarks} de natureza sintética, o Needle-in-a-Haystack e o RULER. O Needle-in-a-Haystack insere um fato isolado, a agulha, em um volume extenso de texto de preenchimento e verifica se o modelo consegue recuperá-lo, variando a posição da agulha e o comprimento do contexto \cite{kamradt2023needle}. O RULER estende essa ideia com múltiplas agulhas, rastreamento de variáveis e tarefas de agregação, mas mantém a construção sintética e controlada \cite{hsieh2024ruler}. Em ambos, a tarefa se reduz a localizar e copiar um trecho específico do contexto, sem exigir que o modelo combine ou sintetize informação proveniente de partes distintas da entrada.

Essa limitação aparece nos próprios resultados da fase preliminar. No GPT-OSS 120B (Tabela~\ref{tab:acuracia}), todas as cinco estratégias alcançam entre 0.87 e 1.00 no Needle-in-a-Haystack, um patamar próximo do teto, enquanto a acurácia no LongBench cai para a faixa de 0.00 a 0.17 e a acurácia no RULER varia de 0.06 a 0.64 conforme a estratégia. A \textit{Semantic Compression}, por exemplo, mantém 1.00 no Needle-in-a-Haystack nesse modelo, mas obtém apenas 0.06 no RULER. Esse contraste indica que o Needle-in-a-Haystack mede sobretudo a capacidade de recuperação posicional, e não a compreensão de contexto longo como um todo. Pontuações próximas do máximo nesse \textit{benchmark}, portanto, não se traduzem necessariamente em bom desempenho em tarefas que exigem integrar ou agregar informação distribuída ao longo do contexto, como as reunidas na suíte revisada.

Como o objetivo desta dissertação é medir quanto de informação uma estratégia de compressão efetivamente preserva, a fase local adota a suíte de \textit{benchmarks} de tarefas realistas descrita na Seção~\ref{subsec:desenho}, mais discriminante para esse propósito do que os \textit{benchmarks} sintéticos empregados na fase preliminar.

A vantagem dessa suíte é que as tarefas dependem de agregação, de raciocínio de múltiplos saltos e da preservação fiel de informação distribuída, que são exatamente as propriedades que a alocação de compressão por partição precisa proteger. Uma estratégia pode obter pontuação alta no Needle-in-a-Haystack ao manter um único trecho, mas falhar em resposta a perguntas de múltiplos saltos se descartar um fato de ligação. Por essa razão, a suíte adotada na fase local mede de forma mais fiel o compromisso entre redução de \textit{tokens} e informação preservada, que é o critério central para a avaliação do \textit{framework} proposto.

\subsection{Resultados da Matriz Local Ampliada}
\label{subsec:resultados}

A matriz local abrange cinco modelos, o Llama 3.1 8B, o Gemma 4 26B, o Qwen 3 30B-A3B, o DeepSeek-R1 32B e o GPT-OSS 20B, e foi construída nas duas etapas descritas na Seção~\ref{subsec:desenho}. A etapa inicial avaliou \textit{Raw}, \textit{Sliding Window}, \textit{Parallel Window}, \textit{Semantic Compression} e RIG em todas as combinações de modelo e orçamento. Posteriormente, LLMLingua-GPT2, LLMLingua-2, Selective Context e CPC-MiniLM foram incorporados como técnicas adicionais, elevando para nove o número de linhas comparáveis em cada bloco da matriz. Essa ampliação foi executada nos cinco modelos e nos dois orçamentos.

A Tabela~\ref{tab:acuracia-atual-8k} apresenta a acurácia por técnica e \textit{benchmark} para o orçamento de 8192 \textit{tokens}, e a Tabela~\ref{tab:acuracia-atual-32k} apresenta os mesmos resultados para o orçamento de 32768 \textit{tokens}. Em ambas, cada técnica ocupa uma linha dentro do respectivo bloco de modelo, e a coluna Média corresponde à média macro das pontuações nos sete \textit{benchmarks} da suíte revisada. Na combinação CPC-MiniLM, DeepSeek-R1 32B e 8192 \textit{tokens}, a pontuação do QMSum e a média macro são calculadas sobre as duas respostas válidas, em vez das três previstas, e identificadas por um asterisco.

{
\footnotesize
\setlength{\tabcolsep}{3pt}
\begin{longtable}{lcc|cc|cc|c|c}
\caption{Acurácia por técnica e \textit{benchmark} nos cinco modelos, orçamento de 8192 \textit{tokens}. As colunas abreviam LongBench (LB), ZeroSCROLLS (ZS), Natural Questions (NQ), TriviaQA (TQA), HotpotQA (HQA), MuSiQue (MSQ) e QMSum (QMS). O asterisco identifica a combinação com uma resposta excluída por \textit{timeout}}
\label{tab:acuracia-atual-8k} \\
\toprule
& \multicolumn{2}{c|}{\textbf{Suítes}} & \multicolumn{2}{c|}{\textbf{Open-domain QA}} & \multicolumn{2}{c|}{\textbf{Multi-hop QA}} & \multicolumn{1}{c|}{\textbf{Sum.}} & \\
\cmidrule(lr){2-3}\cmidrule(lr){4-5}\cmidrule(lr){6-7}\cmidrule(lr){8-8}
\textbf{Técnica} & \textbf{LB} & \textbf{ZS} & \textbf{NQ} & \textbf{TQA} & \textbf{HQA} & \textbf{MSQ} & \textbf{QMS} & \textbf{Média} \\
\midrule
\endfirsthead
\multicolumn{9}{c}{\tablename\ \thetable{} -- continuação} \\
\toprule
& \multicolumn{2}{c|}{\textbf{Suítes}} & \multicolumn{2}{c|}{\textbf{Open-domain QA}} & \multicolumn{2}{c|}{\textbf{Multi-hop QA}} & \multicolumn{1}{c|}{\textbf{Sum.}} & \\
\cmidrule(lr){2-3}\cmidrule(lr){4-5}\cmidrule(lr){6-7}\cmidrule(lr){8-8}
\textbf{Técnica} & \textbf{LB} & \textbf{ZS} & \textbf{NQ} & \textbf{TQA} & \textbf{HQA} & \textbf{MSQ} & \textbf{QMS} & \textbf{Média} \\
\midrule
\endhead
\midrule
\multicolumn{9}{r}{Continua na próxima página} \\
\endfoot
\bottomrule
\endlastfoot
\multicolumn{9}{l}{\textit{Llama 3.1 8B}} \\*
\textit{Raw}                   & \textbf{1.000} & 0.443 & \underline{0.362} & 1.000 & \textbf{0.667} & \textbf{0.417} & 0.273 & 0.594 \\*
\textit{Sliding Window}        & \textbf{1.000} & \textbf{0.642} & \underline{0.362} & 1.000 & \textbf{0.667} & \underline{0.333} & \textbf{0.379} & \textbf{0.626} \\*
\textit{Parallel Window}       & \textbf{1.000} & 0.329 & \underline{0.362} & 1.000 & \textbf{0.667} & \textbf{0.417} & 0.240 & 0.573 \\*
\textit{Semantic Compression}  & \underline{0.938} & 0.377 & \textbf{0.691} & 1.000 & \textbf{0.667} & \underline{0.333} & 0.239 & \underline{0.606} \\*
RIG                            & \textbf{1.000} & 0.404 & \underline{0.362} & 1.000 & \textbf{0.667} & \textbf{0.417} & 0.083 & 0.562 \\*
LLMLingua-GPT2                 & 0.688 & 0.273 & 0.050 & 1.000 & \underline{0.333} & 0.083 & 0.222 & 0.379 \\*
LLMLingua-2                    & 0.812 & \underline{0.619} & 0.080 & 1.000 & \underline{0.333} & \underline{0.333} & 0.262 & 0.492 \\*
Selective Context              & 0.688 & 0.572 & \underline{0.362} & 1.000 & 0.000 & \textbf{0.417} & 0.321 & 0.480 \\*
CPC-MiniLM                     & 0.688 & 0.583 & \underline{0.362} & 1.000 & \underline{0.333} & \textbf{0.417} & \underline{0.335} & 0.531 \\
\midrule
\multicolumn{9}{l}{\textit{Gemma 4 26B}} \\*
\textit{Raw}                   & \textbf{0.938} & \underline{0.333} & \textbf{0.671} & 1.000 & \textbf{0.667} & \textbf{0.333} & 0.000 & \underline{0.563} \\*
\textit{Sliding Window}        & \textbf{0.938} & \underline{0.333} & \textbf{0.671} & 1.000 & \textbf{0.667} & \textbf{0.333} & 0.068 & \textbf{0.573} \\*
\textit{Parallel Window}       & 0.750 & \underline{0.333} & \textbf{0.671} & 1.000 & \textbf{0.667} & \textbf{0.333} & 0.014 & 0.538 \\*
\textit{Semantic Compression}  & \underline{0.812} & \underline{0.333} & 0.000 & 1.000 & \underline{0.333} & \textbf{0.333} & 0.061 & 0.411 \\*
RIG                            & 0.500 & \underline{0.333} & \textbf{0.671} & 1.000 & 0.000 & \textbf{0.333} & 0.088 & 0.418 \\*
LLMLingua-GPT2                 & 0.250 & 0.000 & 0.337 & 1.000 & \underline{0.333} & 0.000 & 0.000 & 0.274 \\*
LLMLingua-2                    & \textbf{0.938} & \textbf{0.583} & \underline{0.341} & 1.000 & \underline{0.333} & 0.000 & 0.007 & 0.457 \\*
Selective Context              & 0.688 & \underline{0.333} & 0.337 & 0.667 & 0.000 & \textbf{0.333} & \underline{0.116} & 0.353 \\*
CPC-MiniLM                     & 0.688 & \underline{0.333} & \textbf{0.671} & 1.000 & \textbf{0.667} & \textbf{0.333} & \textbf{0.138} & 0.547 \\
\midrule
\multicolumn{9}{l}{\textit{Qwen 3 30B-A3B}} \\*
\textit{Raw}                   & \textbf{0.750} & 0.389 & \underline{0.667} & 1.000 & \textbf{0.333} & \textbf{0.333} & 0.141 & \underline{0.516} \\*
\textit{Sliding Window}        & \textbf{0.750} & 0.222 & \underline{0.667} & 1.000 & \underline{0.000} & \textbf{0.333} & 0.000 & 0.425 \\*
\textit{Parallel Window}       & \textbf{0.750} & 0.222 & \textbf{1.000} & 1.000 & \underline{0.000} & \textbf{0.333} & 0.010 & 0.474 \\*
\textit{Semantic Compression}  & \underline{0.625} & 0.417 & 0.333 & 1.000 & \underline{0.000} & \underline{0.000} & \underline{0.153} & 0.361 \\*
RIG                            & \textbf{0.750} & 0.409 & \textbf{1.000} & 1.000 & \underline{0.000} & \textbf{0.333} & \textbf{0.157} & \textbf{0.521} \\*
LLMLingua-GPT2                 & 0.000 & 0.111 & \underline{0.667} & 0.333 & \textbf{0.333} & \underline{0.000} & 0.000 & 0.206 \\*
LLMLingua-2                    & \textbf{0.750} & \textbf{0.592} & \underline{0.667} & 1.000 & \underline{0.000} & \underline{0.000} & 0.000 & 0.430 \\*
Selective Context              & 0.500 & 0.222 & 0.333 & 0.333 & \underline{0.000} & \textbf{0.333} & 0.083 & 0.258 \\*
CPC-MiniLM                     & 0.500 & \underline{0.583} & \underline{0.667} & 0.667 & \underline{0.000} & \textbf{0.333} & 0.005 & 0.394 \\
\midrule
\multicolumn{9}{l}{\textit{DeepSeek-R1 32B}} \\*
\textit{Raw}                   & \textbf{0.750} & 0.620 & 0.164 & 1.000 & \textbf{1.000} & \textbf{0.833} & 0.344 & \textbf{0.673} \\*
\textit{Sliding Window}        & \textbf{0.750} & 0.422 & 0.142 & 1.000 & \textbf{1.000} & 0.500 & 0.335 & 0.593 \\*
\textit{Parallel Window}       & \underline{0.500} & 0.500 & 0.142 & 1.000 & \underline{0.667} & 0.583 & 0.230 & 0.517 \\*
\textit{Semantic Compression}  & \textbf{0.750} & 0.629 & \underline{0.172} & 1.000 & \underline{0.667} & 0.333 & 0.213 & 0.538 \\*
RIG                            & \textbf{0.750} & 0.520 & 0.142 & 1.000 & \underline{0.667} & \underline{0.667} & \underline{0.380} & 0.589 \\*
LLMLingua-GPT2                 & \underline{0.500} & 0.306 & 0.142 & 0.667 & 0.333 & 0.000 & 0.302 & 0.321 \\*
LLMLingua-2                    & \underline{0.500} & 0.650 & 0.150 & 1.000 & \textbf{1.000} & 0.333 & \textbf{0.432} & 0.581 \\*
Selective Context              & \underline{0.500} & \underline{0.690} & \textbf{0.209} & 0.667 & \textbf{1.000} & \underline{0.667} & 0.302 & 0.576 \\*
CPC-MiniLM\textsuperscript{*}  & \underline{0.500} & \textbf{0.743} & \underline{0.172} & 1.000 & \textbf{1.000} & 0.583 & 0.291 & \underline{0.613} \\
\midrule
\multicolumn{9}{l}{\textit{GPT-OSS 20B}} \\*
\textit{Raw}                   & 0.688 & 0.333 & \underline{0.050} & 1.000 & 0.667 & \textbf{0.917} & 0.052 & 0.530 \\*
\textit{Sliding Window}        & \textbf{0.938} & \underline{0.546} & \underline{0.050} & 1.000 & \textbf{1.000} & 0.417 & 0.250 & \textbf{0.600} \\*
\textit{Parallel Window}       & 0.688 & 0.523 & \underline{0.050} & 1.000 & 0.667 & 0.333 & 0.135 & 0.485 \\*
\textit{Semantic Compression}  & \underline{0.750} & 0.476 & 0.006 & 1.000 & 0.500 & 0.000 & \textbf{0.339} & 0.439 \\*
RIG                            & \textbf{0.938} & 0.298 & \underline{0.050} & 1.000 & \underline{0.833} & \underline{0.667} & 0.313 & \underline{0.586} \\*
LLMLingua-GPT2                 & 0.438 & 0.000 & \textbf{0.114} & 1.000 & 0.333 & 0.000 & 0.144 & 0.290 \\*
LLMLingua-2                    & \textbf{0.938} & 0.467 & 0.025 & 1.000 & 0.333 & 0.333 & 0.313 & 0.487 \\*
Selective Context              & 0.312 & 0.391 & 0.028 & 1.000 & 0.667 & 0.333 & 0.255 & 0.427 \\*
CPC-MiniLM                     & 0.688 & \textbf{0.707} & \underline{0.050} & 1.000 & 0.667 & \underline{0.667} & \underline{0.315} & 0.585 \\
\end{longtable}
}

{
\footnotesize
\setlength{\tabcolsep}{3pt}
\begin{longtable}{lcc|cc|cc|c|c}
\caption{Acurácia por técnica e \textit{benchmark} nos cinco modelos, orçamento de 32768 \textit{tokens}. As abreviações de coluna seguem a Tabela~\ref{tab:acuracia-atual-8k}}
\label{tab:acuracia-atual-32k} \\
\toprule
& \multicolumn{2}{c|}{\textbf{Suítes}} & \multicolumn{2}{c|}{\textbf{Open-domain QA}} & \multicolumn{2}{c|}{\textbf{Multi-hop QA}} & \multicolumn{1}{c|}{\textbf{Sum.}} & \\
\cmidrule(lr){2-3}\cmidrule(lr){4-5}\cmidrule(lr){6-7}\cmidrule(lr){8-8}
\textbf{Técnica} & \textbf{LB} & \textbf{ZS} & \textbf{NQ} & \textbf{TQA} & \textbf{HQA} & \textbf{MSQ} & \textbf{QMS} & \textbf{Média} \\
\midrule
\endfirsthead
\multicolumn{9}{c}{\tablename\ \thetable{} -- continuação} \\
\toprule
& \multicolumn{2}{c|}{\textbf{Suítes}} & \multicolumn{2}{c|}{\textbf{Open-domain QA}} & \multicolumn{2}{c|}{\textbf{Multi-hop QA}} & \multicolumn{1}{c|}{\textbf{Sum.}} & \\
\cmidrule(lr){2-3}\cmidrule(lr){4-5}\cmidrule(lr){6-7}\cmidrule(lr){8-8}
\textbf{Técnica} & \textbf{LB} & \textbf{ZS} & \textbf{NQ} & \textbf{TQA} & \textbf{HQA} & \textbf{MSQ} & \textbf{QMS} & \textbf{Média} \\
\midrule
\endhead
\midrule
\multicolumn{9}{r}{Continua na próxima página} \\
\endfoot
\bottomrule
\endlastfoot
\multicolumn{9}{l}{\textit{Llama 3.1 8B}} \\*
\textit{Raw}                   & \textbf{1.000} & 0.619 & \underline{0.362} & 1.000 & \textbf{0.667} & \textbf{0.417} & 0.326 & \textbf{0.627} \\*
\textit{Sliding Window}        & \textbf{1.000} & \underline{0.642} & \underline{0.362} & 1.000 & \textbf{0.667} & \underline{0.333} & \textbf{0.379} & \underline{0.626} \\*
\textit{Parallel Window}       & \textbf{1.000} & 0.329 & \underline{0.362} & 1.000 & \textbf{0.667} & \textbf{0.417} & 0.247 & 0.574 \\*
\textit{Semantic Compression}  & \underline{0.938} & 0.354 & \textbf{0.691} & 1.000 & \textbf{0.667} & \underline{0.333} & 0.158 & 0.592 \\*
RIG                            & \textbf{1.000} & 0.404 & \underline{0.362} & 1.000 & \textbf{0.667} & \textbf{0.417} & 0.083 & 0.562 \\*
LLMLingua-GPT2                 & 0.688 & 0.190 & 0.050 & 1.000 & \underline{0.333} & 0.083 & 0.222 & 0.367 \\*
LLMLingua-2                    & 0.812 & 0.606 & 0.080 & 1.000 & \underline{0.333} & \underline{0.333} & \underline{0.356} & 0.503 \\*
Selective Context              & 0.688 & \textbf{0.656} & \underline{0.362} & 1.000 & 0.000 & \textbf{0.417} & 0.290 & 0.487 \\*
CPC-MiniLM                     & 0.688 & 0.583 & \underline{0.362} & 1.000 & \underline{0.333} & \textbf{0.417} & 0.335 & 0.531 \\
\midrule
\multicolumn{9}{l}{\textit{Gemma 4 26B}} \\*
\textit{Raw}                   & \textbf{0.938} & \textbf{0.583} & \textbf{0.671} & 1.000 & \underline{0.667} & \underline{0.000} & \textbf{0.129} & \underline{0.570} \\*
\textit{Sliding Window}        & \textbf{0.938} & \underline{0.373} & \textbf{0.671} & 1.000 & 0.333 & \textbf{0.333} & 0.000 & 0.521 \\*
\textit{Parallel Window}       & 0.750 & 0.333 & \underline{0.339} & 1.000 & \textbf{1.000} & \textbf{0.333} & 0.014 & 0.538 \\*
\textit{Semantic Compression}  & \underline{0.812} & \textbf{0.583} & 0.000 & 1.000 & 0.333 & \textbf{0.333} & \underline{0.109} & 0.453 \\*
RIG                            & 0.750 & 0.333 & \textbf{0.671} & 1.000 & 0.333 & \textbf{0.333} & 0.054 & 0.496 \\*
LLMLingua-GPT2                 & 0.375 & 0.000 & 0.337 & 1.000 & 0.333 & \underline{0.000} & 0.000 & 0.292 \\*
LLMLingua-2                    & \textbf{0.938} & 0.333 & 0.333 & 1.000 & \underline{0.667} & \underline{0.000} & 0.000 & 0.467 \\*
Selective Context              & 0.688 & 0.222 & 0.337 & 0.667 & 0.000 & \textbf{0.333} & 0.102 & 0.336 \\*
CPC-MiniLM                     & 0.688 & \textbf{0.583} & \textbf{0.671} & 1.000 & \underline{0.667} & \textbf{0.333} & 0.068 & \textbf{0.573} \\
\midrule
\multicolumn{9}{l}{\textit{Qwen 3 30B-A3B}} \\*
\textit{Raw}                   & \textbf{0.750} & \underline{0.472} & \underline{0.667} & 1.000 & \textbf{0.333} & \textbf{0.333} & 0.000 & \underline{0.508} \\*
\textit{Sliding Window}        & \textbf{0.750} & 0.222 & \underline{0.667} & 1.000 & \underline{0.000} & \textbf{0.333} & 0.000 & 0.425 \\*
\textit{Parallel Window}       & \textbf{0.750} & 0.222 & \textbf{1.000} & 1.000 & \underline{0.000} & \textbf{0.333} & 0.010 & 0.474 \\*
\textit{Semantic Compression}  & \underline{0.625} & 0.306 & 0.333 & 1.000 & \underline{0.000} & \underline{0.000} & \underline{0.112} & 0.339 \\*
RIG                            & \textbf{0.750} & 0.409 & \textbf{1.000} & 1.000 & \underline{0.000} & \textbf{0.333} & \textbf{0.157} & \textbf{0.521} \\*
LLMLingua-GPT2                 & 0.000 & 0.000 & \underline{0.667} & 0.333 & \textbf{0.333} & \underline{0.000} & 0.000 & 0.190 \\*
LLMLingua-2                    & \textbf{0.750} & \underline{0.472} & \underline{0.667} & 1.000 & \underline{0.000} & \underline{0.000} & 0.000 & 0.413 \\*
Selective Context              & 0.500 & \underline{0.472} & 0.333 & 0.333 & \underline{0.000} & \textbf{0.333} & 0.000 & 0.282 \\*
CPC-MiniLM                     & 0.500 & \textbf{0.583} & \underline{0.667} & 0.667 & \underline{0.000} & \textbf{0.333} & 0.005 & 0.394 \\
\midrule
\multicolumn{9}{l}{\textit{DeepSeek-R1 32B}} \\*
\textit{Raw}                   & \textbf{0.750} & \underline{0.813} & 0.164 & 1.000 & \textbf{1.000} & \textbf{0.833} & \underline{0.395} & \textbf{0.708} \\*
\textit{Sliding Window}        & \textbf{0.750} & 0.422 & 0.142 & 1.000 & \textbf{1.000} & 0.500 & 0.335 & 0.593 \\*
\textit{Parallel Window}       & \underline{0.500} & 0.500 & 0.142 & 1.000 & \underline{0.667} & 0.583 & 0.269 & 0.523 \\*
\textit{Semantic Compression}  & \textbf{0.750} & 0.624 & \underline{0.172} & 1.000 & \underline{0.667} & 0.333 & 0.307 & 0.550 \\*
RIG                            & \textbf{0.750} & 0.520 & 0.142 & 1.000 & \underline{0.667} & \underline{0.667} & 0.380 & 0.589 \\*
LLMLingua-GPT2                 & \underline{0.500} & 0.389 & 0.142 & 0.667 & 0.333 & 0.000 & 0.302 & 0.333 \\*
LLMLingua-2                    & \underline{0.500} & 0.659 & 0.150 & 1.000 & \textbf{1.000} & 0.333 & \textbf{0.407} & 0.578 \\*
Selective Context              & \underline{0.500} & 0.694 & \textbf{0.209} & 0.667 & \textbf{1.000} & \underline{0.667} & 0.355 & 0.585 \\*
CPC-MiniLM                     & \underline{0.500} & \textbf{0.827} & \underline{0.172} & 1.000 & \textbf{1.000} & 0.583 & 0.354 & \underline{0.634} \\
\midrule
\multicolumn{9}{l}{\textit{GPT-OSS 20B}} \\*
\textit{Raw}                   & 0.688 & \textbf{0.809} & \underline{0.050} & 1.000 & \underline{0.667} & \textbf{0.917} & \textbf{0.329} & \underline{0.637} \\*
\textit{Sliding Window}        & \textbf{0.938} & 0.537 & \underline{0.050} & 1.000 & \textbf{1.000} & \underline{0.750} & 0.255 & \textbf{0.647} \\*
\textit{Parallel Window}       & 0.688 & 0.510 & \underline{0.050} & 1.000 & \textbf{1.000} & 0.667 & 0.217 & 0.590 \\*
\textit{Semantic Compression}  & \underline{0.750} & \underline{0.580} & 0.004 & 1.000 & \underline{0.667} & 0.333 & 0.141 & 0.496 \\*
RIG                            & \textbf{0.938} & 0.298 & \underline{0.050} & 1.000 & \underline{0.667} & 0.667 & 0.204 & 0.546 \\*
LLMLingua-GPT2                 & 0.438 & 0.000 & \textbf{0.114} & 1.000 & 0.333 & 0.000 & 0.144 & 0.290 \\*
LLMLingua-2                    & \textbf{0.938} & 0.142 & 0.025 & 1.000 & 0.333 & 0.333 & 0.295 & 0.438 \\*
Selective Context              & 0.312 & 0.333 & 0.028 & 1.000 & \underline{0.667} & 0.333 & 0.163 & 0.405 \\*
CPC-MiniLM                     & 0.688 & 0.373 & \underline{0.050} & 1.000 & \underline{0.667} & 0.667 & \underline{0.315} & 0.537 \\
\end{longtable}
}

A Tabela~\ref{tab:latencia-compressao-atual} resume a latência média de resposta e a taxa de compressão efetiva por estratégia, medida como a fração de \textit{tokens} retida após o tratamento do contexto, em que o valor 1.000 corresponde à ausência de compressão.

{
\small
\setlength{\tabcolsep}{6pt}
\begin{longtable}{lcc|c}
\caption{Latência média (ms) e taxa de compressão efetiva por técnica nos cinco modelos. A taxa retida é a média das taxas nos dois orçamentos, em que 1.000 corresponde à ausência de compressão. O asterisco identifica a média de latência calculada após a exclusão de uma chamada encerrada por \textit{timeout}}
\label{tab:latencia-compressao-atual} \\
\toprule
& \multicolumn{2}{c|}{\textbf{Latência (ms)}} & \\
\cmidrule(lr){2-3}
\textbf{Técnica} & \textbf{8k} & \textbf{32k} & \textbf{Taxa retida} \\
\midrule
\endfirsthead
\multicolumn{4}{c}{\tablename\ \thetable{} -- continuação} \\
\toprule
& \multicolumn{2}{c|}{\textbf{Latência (ms)}} & \\
\cmidrule(lr){2-3}
\textbf{Técnica} & \textbf{8k} & \textbf{32k} & \textbf{Taxa retida} \\
\midrule
\endhead
\midrule
\multicolumn{4}{r}{Continua na próxima página} \\
\endfoot
\bottomrule
\endlastfoot
\multicolumn{4}{l}{\textit{Llama 3.1 8B}} \\*
\textit{Raw}                   & 8228             & 16588            & 1.000 \\*
\textit{Sliding Window}        & 6764             & \underline{6994} & 0.857 \\*
\textit{Parallel Window}       & 6922             & 7077             & 0.856 \\*
\textit{Semantic Compression}  & 35572            & 87078            & 0.367 \\*
RIG                            & \textbf{4366}    & \textbf{4608}    & 0.692 \\*
LLMLingua-GPT2                 & 8451             & 10127            & 0.508 \\*
LLMLingua-2                    & 9484             & 12385            & 0.470 \\*
Selective Context              & 10344            & 14825            & 0.566 \\*
CPC-MiniLM                     & \underline{6736} & 10345            & 0.524 \\
\midrule
\multicolumn{4}{l}{\textit{Gemma 4 26B}} \\*
\textit{Raw}                   & 15818             & 19361             & 1.000 \\*
\textit{Sliding Window}        & 14752             & \underline{13887} & 0.857 \\*
\textit{Parallel Window}       & 14570             & 14204             & 0.856 \\*
\textit{Semantic Compression}  & 306142            & 354957            & 0.328 \\*
RIG                            & \textbf{13247}    & \textbf{12614}    & 0.692 \\*
LLMLingua-GPT2                 & 18237             & 18864             & 0.508 \\*
LLMLingua-2                    & 17477             & 20048             & 0.470 \\*
Selective Context              & 19508             & 22037             & 0.566 \\*
CPC-MiniLM                     & \underline{14307} & 15755             & 0.524 \\
\midrule
\multicolumn{4}{l}{\textit{Qwen 3 30B-A3B}} \\*
\textit{Raw}                   & 13461             & 20345             & 1.000 \\*
\textit{Sliding Window}        & 12327             & \underline{12340} & 0.857 \\*
\textit{Parallel Window}       & 12371             & 12405             & 0.856 \\*
\textit{Semantic Compression}  & 43051             & 51182             & 0.417 \\*
RIG                            & \textbf{10363}    & \textbf{10381}    & 0.692 \\*
LLMLingua-GPT2                 & 14144             & 17128             & 0.508 \\*
LLMLingua-2                    & 14503             & 18423             & 0.470 \\*
Selective Context              & 16140             & 21096             & 0.566 \\*
CPC-MiniLM                     & \underline{11840} & 15836             & 0.524 \\
\midrule
\multicolumn{4}{l}{\textit{DeepSeek-R1 32B}} \\*
\textit{Raw}                   & 81871             & 111013            & 1.000 \\*
\textit{Sliding Window}        & 76534             & \underline{79862} & 0.857 \\*
\textit{Parallel Window}       & 77592             & 82807             & 0.856 \\*
\textit{Semantic Compression}  & 198572            & 265406            & 0.376 \\*
RIG                            & \textbf{62385}    & \textbf{63864}    & 0.692 \\*
LLMLingua-GPT2                 & \underline{72195} & 80249             & 0.507 \\*
LLMLingua-2                    & 78809             & 90642             & 0.469 \\*
Selective Context              & 108964            & 89828             & 0.566 \\*
CPC-MiniLM                     & 75366\textsuperscript{*} & 83529      & 0.524 \\
\midrule
\multicolumn{4}{l}{\textit{GPT-OSS 20B}} \\*
\textit{Raw}                   & 11233            & 14567             & 1.000 \\*
\textit{Sliding Window}        & \underline{9551} & \textbf{10617}    & 0.857 \\*
\textit{Parallel Window}       & 10150            & 10650             & 0.856 \\*
\textit{Semantic Compression}  & 139627           & 231207            & 0.355 \\*
RIG                            & \textbf{8615}    & \underline{10639} & 0.692 \\*
LLMLingua-GPT2                 & 15449            & 16308             & 0.507 \\*
LLMLingua-2                    & 14546            & 15648             & 0.469 \\*
Selective Context              & 15940            & 17707             & 0.566 \\*
CPC-MiniLM                     & \underline{9214} & 11079             & 0.524 \\
\end{longtable}
}

A leitura conjunta das três tabelas mostra que nenhuma estratégia domina os cinco modelos. Sob o orçamento de 8192 \textit{tokens}, a maior acurácia média é obtida pela \textit{Sliding Window} no Llama 3.1 8B, no Gemma 4 26B e no GPT-OSS 20B, pela RIG no Qwen 3 30B-A3B e pela estratégia \textit{Raw} no DeepSeek-R1 32B. Sob o orçamento de 32768 \textit{tokens}, a estratégia \textit{Raw} lidera no Llama 3.1 8B e no DeepSeek-R1 32B, a RIG permanece à frente no Qwen 3 30B-A3B e a \textit{Sliding Window} lidera no GPT-OSS 20B. No Gemma 4 26B, CPC-MiniLM alcança 0.573 e supera, por pequena margem, a estratégia \textit{Raw}, com 0.570. Essa é a única combinação em que um dos quatro compressores acrescentados ocupa a primeira posição entre as nove técnicas. No DeepSeek-R1 32B, CPC-MiniLM ocupa a segunda posição nos dois orçamentos, com 0.613 e 0.634, mas permanece abaixo da estratégia \textit{Raw}, com 0.673 e 0.708. No GPT-OSS 20B, CPC-MiniLM obtém 0.585 e 0.537, abaixo da \textit{Sliding Window}, que alcança 0.600 e 0.647.

Em latência, a RIG é a estratégia mais rápida entre as que tratam o contexto em todos os cinco modelos sob o orçamento de 8192 \textit{tokens} e em quatro dos cinco sob o orçamento de 32768 \textit{tokens}, com a única exceção do GPT-OSS 20B nesse orçamento, em que a RIG e a \textit{Sliding Window} praticamente empatam. Entre os quatro compressores adicionais, CPC-MiniLM apresenta a menor latência em sete das dez combinações; LLMLingua-GPT2 é mais rápido no Llama 3.1 8B sob 32768 \textit{tokens} e nos dois orçamentos do DeepSeek-R1 32B. A \textit{Semantic Compression} apresenta os maiores tempos em todos os modelos, com valores de uma ordem de grandeza acima das demais e chegando a 354957 ms no Gemma 4 26B sob o orçamento de 32768 \textit{tokens}, uma vez que a compressão semântica exige chamadas adicionais ao próprio modelo. Em taxa de compressão, a \textit{Semantic Compression} é a mais agressiva, retendo entre 0.328 e 0.417 dos \textit{tokens} conforme o modelo. Entre os quatro compressores adicionais, LLMLingua-2 retém a menor fração média, enquanto Selective Context retém a maior.


\subsection{Seleção da Estratégia Base}
\label{subsec:discussao-empirica}

Na etapa inicial da matriz local, a \textit{Sliding Window} apresenta, na média dos dois orçamentos, a maior acurácia entre as estratégias que tratam o contexto no Llama 3.1 8B, no Gemma 4 26B, no DeepSeek-R1 32B e no GPT-OSS 20B. A RIG é superior apenas no Qwen 3 30B-A3B. Consideradas separadamente as dez combinações de modelo e orçamento, a \textit{Sliding Window} alcança a maior acurácia dessa classe em sete, a RIG em duas e a \textit{Parallel Window} em uma. A RIG oferece menor latência média, mas sua acurácia é inferior à da \textit{Sliding Window} em quatro dos cinco modelos. No GPT-OSS 20B, por exemplo, a \textit{Sliding Window} obtém 0.600 e 0.647 nos orçamentos de 8192 e 32768 \textit{tokens}, enquanto a RIG obtém 0.586 e 0.546. Esse resultado, obtido antes da ampliação com os quatro compressores, sustenta a \textit{Sliding Window} como estratégia base da formulação. A etapa posterior não substitui essa comparação; ela caracteriza compressores adicionais que podem compor as opções do MCKP.

A seção seguinte formaliza uma alternativa que distribui a compressão de forma seletiva entre as partições do \textit{prompt}. Embora a \textit{Sliding Window} seja adotada como estratégia base, a formulação admite outras formas de segmentação como ponto de partida para o particionamento.

\section{\textit{Framework} Proposto}
\label{sec:framework}

\subsection{Motivação e Visão Geral}
\label{subsec:motivacao-framework}

O \textit{framework} proposto trata a compressão de \textit{prompt} como um problema de alocação de orçamento. Em vez de aplicar um único compressor e uma única taxa a todo o \textit{prompt}, o \textit{framework} particiona a entrada, estima a importância de cada partição para a tarefa e escolhe, para cada partição, uma opção de compressão que maximiza a informação preservada sob uma restrição global de orçamento de \textit{tokens}. A formulação corresponde a um problema da mochila de múltipla escolha (\textit{Multiple-Choice Knapsack Problem}, MCKP), em que cada partição contribui com exatamente uma opção escolhida entre um conjunto de compressores candidatos.

A Figura~\ref{fig:framework} apresenta a arquitetura conceitual do \textit{framework}, organizada em cinco etapas encadeadas. A figura distingue os artefatos produzidos em cada etapa, desde as partições do \textit{prompt} até a entrada comprimida enviada ao modelo.

% Arquitetura conceitual do framework proposto de compressao seletiva.
% Este diagrama representa a formulacao da dissertacao, nao uma
% implementacao de software ja concluida.

\begin{figure}[htbp]
\centering
\resizebox{0.98\textwidth}{!}{%
\begin{tikzpicture}[
    node distance=0.75cm and 0.9cm,
    stage/.style={rectangle, rounded corners=2pt, draw=dark, thick,
        minimum width=3.25cm, minimum height=1.25cm, align=center,
        font=\small, fill=white},
    artifact/.style={rectangle, draw=dark!65, minimum width=2.8cm,
        minimum height=0.8cm, align=center, font=\scriptsize, fill=light},
    detail/.style={rectangle, draw=dark!45, minimum width=2.45cm,
        minimum height=0.65cm, align=center, font=\scriptsize, fill=white},
    arrow/.style={-{Stealth[length=2.4mm]}, thick, draw=dark},
    dataarrow/.style={-{Stealth[length=2.2mm]}, thick, dashed, draw=dark!70}
]

% Entrada e analise do prompt
\node[artifact, fill=blue!10] (input) {\textit{Prompt} $C$, consulta $q$\\e orcamento $B$};
\node[stage, right=of input, fill=blue!10] (partition) {1. Particionamento\\semantico};
\node[artifact, right=of partition] (parts) {$T=(t_1,\ldots,t_n)$\\particoes ordenadas};
\node[stage, right=of parts, fill=blue!10] (importance) {2. Estimativa de\\importancia $I(t_j;q)$};

\draw[arrow] (input) -- (partition);
\draw[arrow] (partition) -- (parts);
\draw[arrow] (parts) -- (importance);

% Componentes da importancia
\node[detail, below=0.55cm of importance, xshift=-2.7cm] (relevance) {Relevancia para\\a tarefa};
\node[detail, right=0.25cm of relevance] (density) {Densidade de\\informacao};
\node[detail, right=0.25cm of density] (position) {Saliencia\\posicional};

\draw[dataarrow] (relevance.north) |- (importance.south);
\draw[dataarrow] (density.north) -- (importance.south);
\draw[dataarrow] (position.north) |- (importance.south);

% Geracao de opcoes
\node[stage, below=2.35cm of parts, fill=green!10] (options) {3. Geracao de opcoes\\por particao};
\node[artifact, left=of options, minimum width=3.25cm] (pool) {Compressores \textit{text-in/text-out}\\identidade e candidatos};
\node[artifact, right=of options, minimum width=3.25cm] (optiondata) {$O_j=\{o_{j,0},\ldots,o_{j,m}\}$\\custo $c_{j,o}$ e qualidade $Q_{j,o}$};

\draw[arrow] (parts.south) -- (options.north);
\draw[dataarrow] (importance.east) -- ++(0.45,0) |- (optiondata.east);
\draw[dataarrow] (pool) -- (options);
\draw[arrow] (options) -- (optiondata);

% Otimizacao e saida
\node[stage, below=1.05cm of optiondata, fill=orange!15, minimum width=4.1cm]
    (solver) {4. Otimizacao MCKP\\por programacao dinamica};
\node[detail, left=of solver, minimum width=3.25cm] (constraints)
    {Uma opcao por particao\\$\sum c_{j,o}\leq B$};
\node[detail, right=of solver, minimum width=3.25cm] (transition)
    {Qualidade preservada\\e penalizacao de transicao};

\draw[arrow] (optiondata) -- (solver);
\draw[dataarrow] (constraints) -- (solver);
\draw[dataarrow] (transition) -- (solver);

\node[artifact, below=1cm of solver, fill=green!8] (compressed)
    {\textit{Prompt} comprimido\\com no maximo $B$ tokens};
\node[stage, left=of compressed, fill=green!10] (rebuild)
    {5. Reconstrucao na\\ordem original};
\node[artifact, right=of compressed, fill=purple!10] (llm)
    {Modelo de linguagem\\acessado via API};

\draw[arrow] (solver.south) -- ++(0,-0.35) -| (rebuild.north);
\draw[arrow] (rebuild) -- (compressed);
\draw[arrow] (compressed) -- (llm);

\end{tikzpicture}%
}
\caption{Arquitetura conceitual do \textit{framework} proposto. O sistema transforma cada particao em alternativas com custo e qualidade, seleciona uma alternativa por particao sob o orcamento global e reconstrói o \textit{prompt} antes da chamada ao modelo.}
\label{fig:framework}
\end{figure}


\subsection{Particionamento e Orçamento}
\label{subsec:particionamento}

Seja $C$ o \textit{prompt} de entrada, com comprimento $|C|$ medido em \textit{tokens}. O \textit{framework} particiona $C$ em uma sequência de $n$ partições $T = (t_1, t_2, \ldots, t_n)$, de modo que cada partição corresponda a uma unidade semântica coerente, como uma instrução, uma pergunta, um documento recuperado ou um trecho de fundo. Define-se um orçamento efetivo de contexto $B \leq |C|$, que representa o número máximo de \textit{tokens} que o \textit{prompt} comprimido pode ocupar após a seleção das opções de compressão.

O particionador utilizado na execução principal é estrutural e produz unidades disjuntas, sem sobreposição entre trechos vizinhos. Parágrafos consecutivos são agrupados até o limite de 400 \textit{tokens}. Quando um parágrafo excede esse limite, ele é dividido por fronteiras de sentença; sentenças ainda maiores são separadas por palavras como último recurso. Esse procedimento preserva a cobertura do texto original e impede a duplicação de \textit{tokens}, condição necessária para que os custos das opções sejam aditivos. Uma variante de particionamento semântico permanece implementada, mas não integra a configuração principal. Para o controle uniforme, o contexto completo é representado como uma única partição.

\subsection{Função de Importância}
\label{subsec:importancia}

Para cada partição $t_j$, define-se uma função de importância $I(t_j; q)$ que quantifica a relevância da partição para a consulta $q$. A função combina múltiplos fatores por meio de uma soma ponderada, conforme a Equação~\ref{eq:importancia}.

\begin{equation}
\label{eq:importancia}
I(t_j; q) = \sum_{i} w_i \, \varphi_i(t_j, q),
\end{equation}

onde cada $\varphi_i(t_j, q)$ é um componente de relevância e $w_i$ é o peso associado. Os componentes considerados incluem a relevância da partição para a tarefa, a densidade de informação da partição e a saliência posicional, esta última motivada pelo viés \textit{lost in the middle} discutido na Seção~\ref{subsec:problema-critico}. A função de importância orienta a alocação de orçamento, atribuindo maior prioridade às partições cuja perda de conteúdo teria maior impacto sobre a resposta.

Na implementação atual, a relevância $R(t_j,q)$ combina similaridade semântica e sobreposição lexical. A similaridade semântica corresponde ao cosseno entre os \textit{embeddings} normalizados da partição e da consulta, produzidos pelo all-MiniLM-L6-v2. O componente lexical é calculado por TF-IDF. Os dois sinais são combinados segundo

\begin{equation}
\label{eq:relevancia-implementada}
R(t_j,q)=0.7\,R_{\mathrm{sem}}(t_j,q)+0.3\,R_{\mathrm{lex}}(t_j,q).
\end{equation}

A densidade factual $D(t_j)$ é aproximada pela proporção de ocorrências numéricas, datas e palavras iniciadas por letra maiúscula em relação ao número de palavras da partição. Esse componente constitui um indicador lexical de concentração factual, e não um reconhecedor de entidades. A saliência posicional $P(t_j)$ utiliza uma curva central, com valor maior para partições próximas ao centro da sequência:

\begin{equation}
\label{eq:saliencia-central}
P(t_j)=1-\left|2\frac{j-1}{n-1}-1\right|,
\end{equation}

para $n>1$. Antes da combinação, $R$, $D$ e $P$ são normalizados por \textit{min--max} entre as partições da mesma entrada. Quando todos os valores de um componente são iguais, o componente normalizado recebe valor 1 em todas as partições. Com os pesos da configuração experimental, a função efetivamente utilizada é

\begin{equation}
\label{eq:importancia-implementada}
I(t_j;q)=0.6\,\widehat{R}(t_j,q)
        +0.2\,\widehat{D}(t_j)
        +0.2\,\widehat{P}(t_j),
\end{equation}

onde o acento circunflexo indica o valor normalizado. Os pesos são mantidos fixos na execução principal e tratados como componentes sujeitos a ablação, uma vez que sua combinação representa uma hipótese operacional sobre relevância, densidade e posição.

\subsection{Opções de Compressão e Modelo de Qualidade}
\label{subsec:opcoes}

Para cada partição $t_j$, o \textit{framework} considera um conjunto de opções de compressão $O_j = \{o_{j,0}, o_{j,1}, \ldots, o_{j,m}\}$. Cada opção $o = (A_k, p)$ corresponde à aplicação de um compressor $A_k$ com um parâmetro de operação $p$, como uma taxa de compressão alvo. A opção $o_{j,0}$ corresponde ao compressor identidade $A_0$, que preserva a partição integralmente e serve de referência para o custo e a qualidade máximos daquela partição.

Cada opção possui um custo $c_{j,o}$, medido pelo número de \textit{tokens} da partição após a compressão, e uma qualidade $Q_{j,o}$, que estima a informação preservada pela opção. A qualidade combina a importância da partição com a fidelidade da opção de compressão, conforme a Equação~\ref{eq:qualidade}.

\begin{equation}
\label{eq:qualidade}
Q_{j,o} = I(t_j; q) \cdot f_{j,o},
\end{equation}

onde $f_{j,o} \in [0,1]$ representa a fração de informação da partição preservada pela opção $o$. O compressor identidade satisfaz $f_{j,0} = 1$, e opções mais agressivas apresentam valores menores de $f_{j,o}$, capturando o compromisso entre redução de \textit{tokens} e perda de conteúdo.

Na implementação atual, a fidelidade semântica é calculada por cobertura entre \textit{embeddings}. A partição original e sua representação comprimida são divididas em trechos de 600 caracteres. Para limitar o custo, no máximo 48 trechos de cada texto são mantidos por amostragem uniforme. Seja $S_j$ o conjunto de trechos da partição original e $S'_{j,o}$ o conjunto correspondente à opção comprimida. A fidelidade semântica é

\begin{equation}
\label{eq:fidelidade-semantica-implementada}
f^{\mathrm{sem}}_{j,o}
=
\frac{1}{|S_j|}
\sum_{s\in S_j}
\max_{u\in S'_{j,o}}
\cos\bigl(E(s),E(u)\bigr),
\end{equation}

onde $E(\cdot)$ é o codificador all-MiniLM-L6-v2. A medida estima quanto do conteúdo semântico de cada trecho original encontra uma representação próxima no texto comprimido. Ela não utiliza a resposta esperada do \textit{benchmark}, evitando que informação do conjunto de avaliação seja incorporada à decisão do MCKP.

Como proteção factual adicional, calcula-se a retenção numérica $r^{\mathrm{num}}_{j,o}$, definida como a fração de números, datas, percentuais, valores e horários da partição original que permanecem na opção comprimida. Quando essa retenção é aplicável e fica abaixo do limiar de 0.9, a fidelidade da opção é ajustada por

\begin{equation}
\label{eq:fidelidade-ajustada}
f_{j,o}
=
\min\left(f^{\mathrm{sem}}_{j,o},r^{\mathrm{num}}_{j,o}\right).
\end{equation}

Nos demais casos, $f_{j,o}=f^{\mathrm{sem}}_{j,o}$. A identidade recebe fidelidade 1 e a omissão recebe fidelidade 0. O conjunto principal de opções contém identidade, omissão, CPC-MiniLM, LLMLingua-2 e Selective Context. Os três compressores são materializados com taxas nominais de retenção de 0.5 e 0.3. Se o custo serializado das identidades exceder o orçamento, acrescenta-se a cada compressor uma taxa adaptativa

\begin{equation}
\label{eq:taxa-adaptativa}
\rho_{\mathrm{adapt}}
=
\max\left(0.05,
\min\left(1,\;0.9\frac{B}{C_{\mathrm{id}}}\right)\right),
\end{equation}

onde $C_{\mathrm{id}}$ é o custo total das partições sem compressão. A taxa mínima de 0.05 também é avaliada quando difere da taxa adaptativa. O custo $c_{j,o}$ é medido após a materialização da opção e inclui, de forma conservadora, o separador empregado na reconstrução. Assim, a taxa solicitada ao compressor não substitui a medição do custo efetivamente produzido.

\subsection{Restrições de Compatibilidade}
\label{subsec:compatibilidade}

Nem toda combinação de opções entre partições adjacentes é adequada, pois transições abruptas de estilo ou de nível de compressão entre trechos vizinhos podem prejudicar a coerência do \textit{prompt} reconstruído. Define-se uma função de compatibilidade $M(o, o')$ entre opções de partições adjacentes, que penaliza combinações incompatíveis. A compatibilidade é incorporada como uma penalização de transição na função objetivo, por meio de um termo $\mu \sum_{j} d(o_j, o_{j+1})$, onde $d(\cdot, \cdot)$ mede a distância entre opções adjacentes e $\mu$ controla a intensidade da penalização.

Na configuração principal, a distância depende da família do compressor:

\begin{equation}
\label{eq:distancia-familia}
d(o,o')=
\begin{cases}
0, & \text{se } A(o)=A(o'),\\
1, & \text{caso contrário},
\end{cases}
\end{equation}

onde $A(o)$ identifica o compressor associado à opção. O valor principal é $\mu=0.1$, interpretado como uma regularização equivalente a 10\% da escala unitária de valor para cada troca de família entre partições adjacentes. Esse valor constitui uma escolha inicial do protocolo, e não um parâmetro derivado da formulação. As execuções com $\mu=0$ e $\mu=0.5$ são mantidas como ablações para distinguir o MCKP clássico de uma penalização mais intensa.

\subsection{Formulação como Extensão do Problema da Mochila de Múltipla Escolha}
\label{subsec:mckp}

A seleção das opções de compressão é formulada como uma extensão do MCKP com penalização entre escolhas adjacentes. Para cada partição $t_j$ e cada opção $o \in O_j$, define-se uma variável binária $x_{j,o} \in \{0,1\}$, que assume o valor 1 quando a opção $o$ é escolhida para a partição $t_j$. Denota-se por $o_j$ a única opção cuja variável assume valor 1 na partição $j$. O problema consiste em maximizar a qualidade total preservada sujeita ao orçamento e à escolha de exatamente uma opção por partição, conforme as Equações~\ref{eq:mckp-obj} a \ref{eq:mckp-bin}. Quando $\mu=0$, a formulação se reduz ao MCKP clássico; quando $\mu>0$, o termo de transição incorpora a coerência entre partições vizinhas.

\begin{align}
\max_{x} \quad & \sum_{j=1}^{n} \sum_{o \in O_j} x_{j,o} \, Q_{j,o}
    - \mu \sum_{j=1}^{n-1} d(o_j, o_{j+1}) \label{eq:mckp-obj} \\
\text{s.a.} \quad & \sum_{j=1}^{n} \sum_{o \in O_j} x_{j,o} \, c_{j,o} \leq B \label{eq:mckp-orc} \\
& \sum_{o \in O_j} x_{j,o} = 1, \quad j = 1, \ldots, n \label{eq:mckp-um} \\
& x_{j,o} \in \{0,1\}. \label{eq:mckp-bin}
\end{align}

A Equação~\ref{eq:mckp-obj} maximiza a qualidade total preservada, descontada da penalização de transição entre partições adjacentes. A Equação~\ref{eq:mckp-orc} impõe o orçamento efetivo de contexto. A Equação~\ref{eq:mckp-um} garante que exatamente uma opção seja escolhida por partição, característica que distingue o núcleo MCKP do problema da mochila clássico.

\subsection{Solução por Programação Dinâmica}
\label{subsec:dp}

O MCKP com orçamento inteiro admite solução por programação dinâmica. Como a função objetivo inclui uma penalização entre opções adjacentes, o estado também deve registrar a opção escolhida para a última partição. Define-se $D[j,b,o]$ como a qualidade máxima preservada considerando as $j$ primeiras partições sob um orçamento parcial $b$, quando a opção $o\in O_j$ é escolhida para a partição $t_j$. A recorrência é apresentada na Equação~\ref{eq:dp}.

\begin{equation}
\label{eq:dp}
D[j,b,o] = Q_{j,o} + \max_{o' \in O_{j-1}}
\Big\{D[j-1,b-c_{j,o},o'] - \mu\,d(o',o)\Big\},
\end{equation}

Para a primeira partição, define-se $D[1,b,o]=Q_{1,o}$ quando $c_{1,o}\leq b$, e $D[1,b,o]=-\infty$ caso contrário. A solução ótima corresponde a $\max_{o\in O_n}D[n,B,o]$, e as opções escolhidas são recuperadas por retrocesso (\textit{backtracking}) sobre a tabela de decisão. Para no máximo $m$ opções por partição, a complexidade temporal é $O(nBm^2)$ e a tabela completa ocupa $O(nBm)$ posições; a memória pode ser reduzida quando apenas o valor ótimo, sem o retrocesso, é necessário.

Na implementação de referência, o eixo de orçamento utiliza unidades inteiras de \textit{tokens} (\texttt{budget\_bucket=1}), preservando a solução exata para o modelo de custos declarado. Uma quantização opcional por blocos permanece disponível como aproximação conservadora para instâncias maiores, mas não é utilizada nos experimentos principais. Os custos das opções incluem a serialização necessária à reconstrução, e o contexto final é recontado antes da chamada ao modelo para verificar a restrição de orçamento.

A Figura~\ref{fig:mckp-dp-flow} detalha o fluxo proposto para o preenchimento da tabela de programação dinâmica e para a recuperação das opções selecionadas.

% Fluxo conceitual da solucao do MCKP por programacao dinamica.

\begin{figure}[htbp]
\centering
\resizebox{!}{0.70\textheight}{%
\begin{tikzpicture}[
    node distance=0.7cm,
    startstop/.style={rectangle, rounded corners=3pt, draw=dark, thick,
        minimum width=4.8cm, minimum height=0.85cm, align=center,
        font=\small\bfseries, fill=blue!10},
    process/.style={rectangle, draw=dark, thick, minimum width=5.8cm,
        minimum height=0.9cm, align=center, font=\small, fill=white},
    decision/.style={diamond, aspect=2.4, draw=dark, thick,
        minimum width=3.4cm, minimum height=0.85cm, align=center,
        font=\small, fill=orange!15},
    arrow/.style={-{Stealth[length=2.4mm]}, thick, draw=dark},
    loop/.style={-{Stealth[length=2.2mm]}, thick, dashed, draw=dark!70}
]

\node[startstop] (start) {Entrada: conjuntos $O_1,\ldots,O_n$, orcamento $B$ e penalizacao $\mu$};
\node[process, below=of start] (init) {Inicializar os estados da primeira particao e a tabela de predecessores};
\node[process, below=of init] (part) {Selecionar a proxima particao $j$};
\node[process, below=of part] (budget) {Selecionar o proximo estado $(b,o)$, com $o\in O_j$ e $c_{j,o}\leq b$};
\node[process, below=of budget] (evaluate) {Avaliar cada opcao anterior $o'\in O_{j-1}$};
\node[process, below=of evaluate, fill=orange!10] (update)
    {$D[j,b,o]\leftarrow Q_{j,o}+\max_{o'}\{D[j-1,b-c_{j,o},o']-\mu d(o',o)\}$};
\node[decision, below=0.9cm of update] (moreb) {Restam estados $(b,o)$?};
\node[decision, below=1cm of moreb] (morej) {Restam particoes?};
\node[process, below=1cm of morej, fill=green!10] (backtrack)
    {Retroceder a partir de $\max_o D[n,B,o]$ e recuperar uma opcao por particao};
\node[process, below=of backtrack] (rebuild)
    {Concatenar as representacoes escolhidas na ordem original};
\node[startstop, below=of rebuild, fill=green!15] (end)
    {Saida: \textit{prompt} comprimido sob o orcamento $B$};

\draw[arrow] (start) -- (init);
\draw[arrow] (init) -- (part);
\draw[arrow] (part) -- (budget);
\draw[arrow] (budget) -- (evaluate);
\draw[arrow] (evaluate) -- (update);
\draw[arrow] (update) -- (moreb);
\draw[loop] (moreb.east) -- ++(2.2,0) node[above, font=\scriptsize] {sim} |- (budget.east);
\draw[arrow] (moreb) -- node[right, font=\scriptsize] {nao} (morej);
\draw[loop] (morej.west) -- ++(-2.2,0) node[above, font=\scriptsize] {sim} |- (part.west);
\draw[arrow] (morej) -- node[right, font=\scriptsize] {nao} (backtrack);
\draw[arrow] (backtrack) -- (rebuild);
\draw[arrow] (rebuild) -- (end);

\end{tikzpicture}%
}
\caption{Fluxo proposto para resolver a selecao de compressores por programacao dinamica. A tabela registra a melhor qualidade acumulada para cada particao, orcamento parcial e ultima opcao; o retrocesso recupera as opcoes usadas na reconstrucao.}
\label{fig:mckp-dp-flow}
\end{figure}


\subsection{Algoritmo do \textit{Framework}}
\label{subsec:algoritmo}

O funcionamento do \textit{framework} pode ser resumido em cinco etapas encadeadas.

\begin{enumerate}[leftmargin=1.8em]
\item \textbf{Particionamento}. O \textit{prompt} $C$ é dividido em partições $T = (t_1, \ldots, t_n)$ segundo unidades semânticas coerentes.

\item \textbf{Estimativa de importância}. Para cada partição, calcula-se a importância $I(t_j; q)$ segundo a Equação~\ref{eq:importancia}.

\item \textbf{Geração de opções}. Para cada partição, aplica-se o conjunto de compressores candidatos, produzindo opções com seus custos $c_{j,o}$ e qualidades $Q_{j,o}$.

\item \textbf{Otimização}. Resolve-se a extensão do MCKP das Equações~\ref{eq:mckp-obj} a \ref{eq:mckp-bin} por programação dinâmica, respeitando o orçamento $B$ e as restrições de compatibilidade.

\item \textbf{Reconstrução}. As opções escolhidas são concatenadas na ordem original das partições, produzindo o \textit{prompt} comprimido enviado ao modelo.
\end{enumerate}

A Figura~\ref{fig:sequence-diagram} apresenta a sequência de interação implementada entre os componentes necessários para executar essas etapas.

% Sequencia conceitual entre os componentes do framework proposto.

\begin{figure}[htbp]
\centering
\resizebox{0.98\textwidth}{!}{%
\begin{tikzpicture}[
    actor/.style={rectangle, rounded corners=2pt, minimum width=2.15cm,
        minimum height=0.7cm, draw=dark, thick, fill=blue!10,
        font=\scriptsize\bfseries, align=center},
    lifeline/.style={dashed, draw=dark!55},
    message/.style={-{Stealth[length=2mm]}, thick, draw=dark},
    return/.style={-{Stealth[length=2mm]}, dashed, draw=dark!75},
    note/.style={rectangle, draw=dark!45, fill=orange!10,
        font=\tiny, align=left, inner sep=3pt}
]

\node[actor] (app) at (0,0) {Aplicacao};
\node[actor] (partitioner) at (3.1,0) {Particionador};
\node[actor] (estimator) at (6.2,0) {Estimador de\\importancia};
\node[actor] (pool) at (9.3,0) {Gerador de\\opcoes};
\node[actor] (solver) at (12.4,0) {Otimizador\\MCKP};
\node[actor] (assembler) at (15.5,0) {Reconstrutor};
\node[actor] (llm) at (18.6,0) {LLM via API};

\foreach \x in {app,partitioner,estimator,pool,solver,assembler,llm} {
    \draw[lifeline] (\x.south) -- ++(0,-11.6);
}

\draw[message] ([yshift=-0.65cm]app.south) -- node[above, font=\tiny] {1: comprimir$(C,q,B)$} ([yshift=-0.65cm]partitioner.south);
\draw[return] ([yshift=-1.3cm]partitioner.south) -- node[above, font=\tiny] {$T=(t_1,\ldots,t_n)$} ([yshift=-1.3cm]app.south);

\node[note] at (4.65,-2.1) {Para cada particao $t_j$};
\draw[message] ([yshift=-2.8cm]app.south) -- node[above, font=\tiny] {2: estimar$(t_j,q)$} ([yshift=-2.8cm]estimator.south);
\draw[return] ([yshift=-3.5cm]estimator.south) -- node[above, font=\tiny] {$I(t_j;q)$} ([yshift=-3.5cm]app.south);

\draw[message] ([yshift=-4.25cm]app.south) -- node[above, font=\tiny] {3: gerar$(t_j,I_j)$} ([yshift=-4.25cm]pool.south);
\node[note, right=0.2cm of pool, yshift=-4.9cm] {Aplicar identidade e\\compressores candidatos;\\medir $c_{j,o}$ e $Q_{j,o}$};
\draw[return] ([yshift=-5.7cm]pool.south) -- node[above, font=\tiny] {$O_j$} ([yshift=-5.7cm]app.south);

\draw[message] ([yshift=-6.55cm]app.south) -- node[above, font=\tiny] {4: resolver$(\{O_j\},B,\mu)$} ([yshift=-6.55cm]solver.south);
\node[note, right=0.2cm of solver, yshift=-7.25cm] {Preencher tabela $D$\\e executar retrocesso};
\draw[return] ([yshift=-8cm]solver.south) -- node[above, font=\tiny] {$o_1^*,\ldots,o_n^*$} ([yshift=-8cm]app.south);

\draw[message] ([yshift=-8.75cm]app.south) -- node[above, font=\tiny] {5: reconstruir$(T,O^*)$} ([yshift=-8.75cm]assembler.south);
\draw[return] ([yshift=-9.45cm]assembler.south) -- node[above, font=\tiny] {$C'$ com $|C'|\leq B$} ([yshift=-9.45cm]app.south);

\draw[message] ([yshift=-10.2cm]app.south) -- node[above, font=\tiny] {6: gerar$(C',q)$} ([yshift=-10.2cm]llm.south);
\draw[return] ([yshift=-11cm]llm.south) -- node[above, font=\tiny] {resposta} ([yshift=-11cm]app.south);

\end{tikzpicture}%
}
\caption{Sequencia conceitual da aplicacao do \textit{framework}. Para cada particao, sao calculadas a importancia e as alternativas de compressao; o otimizador seleciona as alternativas globalmente antes da reconstrucao e da chamada ao modelo.}
\label{fig:framework-sequence}
\end{figure}


\subsection{Conjunto de Compressores Compatíveis}
\label{subsec:pool}

O \textit{framework} atua como uma camada de decisão acima de compressores existentes, e não como um novo compressor. Por essa razão, o conjunto de opções por partição é composto por compressores do tipo \textit{text-in}/\textit{text-out}, isto é, que recebem texto e produzem texto comprimido, o que garante compatibilidade com modelos acessados via API. LLMLingua é empregado como baseline externo, pois constitui o ponto de comparação comum aos três trabalhos posteriores revisados na Seção~\ref{subsec:compressao-prompt}. LLMLingua-2, Selective Context e CPC são avaliados como candidatos às opções do MCKP. A Tabela~\ref{tab:compressores} resume os respectivos papéis experimentais.

\begin{table}[htbp]
\centering
\caption{Baseline e compressores candidatos considerados para o MCKP}
\label{tab:compressores}
\begin{tabular}{>{\raggedright\arraybackslash}p{3.2cm}>{\raggedright\arraybackslash}p{2.5cm}>{\raggedright\arraybackslash}p{6.0cm}}
\toprule
\textbf{Compressor} & \textbf{Papel} & \textbf{Característica relevante} \\
\midrule
Identidade ($A_0$) & Referência interna & Preserva a partição integralmente, com fidelidade máxima. \\
\addlinespace
LLMLingua-GPT2 \cite{jiang2023llmlingua} & Baseline com \textit{scorer} adaptado & Implementação do pacote \texttt{llmlingua}, com GPT-2 em CPU para reduzir o custo local. \\
\addlinespace
LLMLingua-2 \cite{pan2024llmlingua2} & Candidato MCKP & Classificação de palavras treinada por destilação, agnóstica à pergunta e com baixa latência. \\
\addlinespace
Selective Context \cite{li2023selective} & Candidato MCKP & Remoção de unidades lexicais por auto-informação, sem treinamento específico. \\
\addlinespace
CPC-MiniLM \cite{li2024cpc} & Aproximação candidata & Seleção de sentenças condicionada à pergunta com \textit{embeddings} do all-MiniLM-L6-v2. \\
\bottomrule
\end{tabular}
\end{table}

A recomendação dos três candidatos se baseia em sua complementaridade. LLMLingua-2 opera em nível de palavra e amortiza o custo de compressão entre consultas, Selective Context introduz uma alternativa sem treinamento específico e CPC-MiniLM preserva sentenças completas com seleção condicionada à pergunta. A escolha por compressores \textit{text-in}/\textit{text-out} mantém a propriedade não invasiva do \textit{framework}, pois nenhuma etapa exige acesso aos parâmetros internos, retreinamento ou modificação arquitetural do modelo base. O conjunto candidato pode ser estendido com outros compressores compatíveis sem alteração da formulação, uma vez que cada compressor é incorporado apenas como uma nova opção por partição. Nas dez combinações de modelo e orçamento, CPC-MiniLM apresenta a maior média entre os quatro compressores em oito, enquanto LLMLingua-2 lidera nas duas combinações do Qwen 3 30B-A3B. Selective Context preserva a maior fidelidade semântica segundo a métrica adotada. Esses resultados sustentam LLMLingua-2, Selective Context e CPC-MiniLM como conjunto inicial de opções do MCKP; LLMLingua-GPT2 permanece como baseline externo. Os resultados são apresentados na Seção~\ref{subsec:resultados-triagem-compressores}.

Na configuração primária definida na Seção~\ref{subsec:selecao-modelo-mckp}, com Llama~3.1 8B e janela de 8.192 tokens, CPC-MiniLM também foi o compressor individual com melhor resultado downstream. Nos 22 casos da matriz exploratória, obteve médias macro e micro de 0,531 e 0,538, respectivamente, com latência média de 6,74~s e retenção de 0,517. LLMLingua-2 atingiu média macro de 0,492 e a menor retenção (0,457), enquanto Selective Context alcançou média macro de 0,480 e a maior fidelidade semântica (0,733). LLMLingua-GPT2 obteve média macro de 0,379. Portanto, CPC-MiniLM é adotado como referência primária quando se exige um único compressor, sem excluir os demais do espaço de decisão do MCKP, pois suas características complementares podem ser úteis para segmentos específicos. Esses valores são tratados como evidência exploratória, uma vez que a matriz contém uma execução por caso e utiliza instâncias sintéticas nos conjuntos derivados do LongBench.

\subsection{Implementação de Referência}
\label{subsec:implementacao}

O repositório contém uma implementação integrada no subdiretório \texttt{mckp}. O módulo \path{partitioner.py} produz partições disjuntas; \path{importance.py} calcula $I(t_j;q)$; \path{options.py} materializa e avalia as opções; \path{solver.py} executa a programação dinâmica e o retrocesso; e \path{reconstructor.py} preserva a ordem das opções selecionadas sem introduzir marcadores não contabilizados. A classe \texttt{MCKPStrategy}, em \path{strategy.py}, coordena essas etapas e expõe a estratégia ao executor de \textit{benchmarks}.

\paragraph{Separação em relação à etapa inicial da matriz local.}
A estratégia denominada \textit{Semantic Compression} na etapa inicial da matriz local não implementava LLMLingua. Em vez disso, ela solicitava que a própria LLM reescrevesse o contexto por meio de uma instrução de compressão, conforme a abordagem de Gilbert et al. \cite{gilbert2023semanticcompression}. O repositório também continha um \texttt{PerplexityCompressor} experimental, que selecionava \textit{tokens} pela perda calculada com GPT-2 e podia ser descrito apenas como inspirado no critério do LLMLingua. Esse protótipo não estava registrado entre as estratégias avaliadas. Para a ampliação da matriz foram criadas entradas explícitas e independentes para LLMLingua, LLMLingua-2, Selective Context e CPC-MiniLM. Assim, resultados da \textit{Semantic Compression} ou do protótipo de perplexidade não são apresentados como resultados de LLMLingua.

\paragraph{Adaptações para execução local.}
Os experimentos utilizam uma LLM alvo executada localmente pelo servidor Ollama 0.32.3, de modo que o compressor e o modelo gerador compartilham os recursos computacionais da mesma máquina. Cada chamada informa explicitamente \texttt{num\_ctx}, reserva até 500 \textit{tokens} para a saída, utiliza temperatura 0 e fixa a semente em 42. O orçamento do MCKP desconta da janela o \textit{template} completo sem contexto, a consulta, a reserva de saída e uma margem de segurança de 256 \textit{tokens}. O contexto reconstruído e o \textit{prompt} final são recontados antes da chamada, enquanto a contagem real reportada pelo Ollama é registrada para avaliar a diferença entre o contador da otimização e o tokenizador do modelo alvo.

O baseline invoca a classe \texttt{PromptCompressor} do pacote \texttt{llmlingua} 0.2.2, com \texttt{use\_llmlingua2=False}, taxa de retenção de 0,5 e execução em CPU. Para reduzir a concorrência de memória com os modelos servidos pelo Ollama, o modelo causal auxiliar é o \texttt{gpt2} \cite{radford2019gpt2}, que já se encontrava no \textit{cache} local e é substancialmente menor que um auxiliar de 7 bilhões de parâmetros. A escolha é compatível com a arquitetura do LLMLingua e com a análise de diferentes modelos auxiliares apresentada por Jiang et al. \cite{jiang2023llmlingua}; entretanto, difere das configurações alinhadas, Alpaca-7B e GPT2-Alpaca, avaliadas no artigo. Por isso, a configuração é reportada como \textbf{LLMLingua-GPT2}, uma vez que o algoritmo provém da implementação do LLMLingua, enquanto o \textit{scorer} não ajustado e o dispositivo são escolhas locais que devem acompanhar a interpretação dos resultados.

LLMLingua-2 também é executado pelo pacote \texttt{llmlingua} 0.2.2, com o \textit{checkpoint} \path{microsoft/llmlingua-2-xlm-roberta-large-meetingbank} e \texttt{device\_map=cpu}. A fixação explícita do dispositivo foi necessária porque a inicialização padrão do pacote pressupunha CUDA, indisponível nessa configuração. A biblioteca Transformers foi fixada na versão 4.46.3, pois o LLMLingua 0.2.2 manipula o formato legado do \textit{cache} de chaves e valores; a versão 5.12.1 instalada inicialmente passou a retornar \texttt{DynamicCache} e provocava falhas durante a compressão iterativa. Selective Context é executado pelo pacote \texttt{selective-context} 0.1.4 com GPT-2, idioma inglês e \texttt{reduce\_ratio=1-rate}, pois sua API recebe a fração removida, ao passo que o executor comum recebe a fração retida. Essa integração utiliza spaCy 3.8.14 e o modelo \texttt{en\_core\_web\_sm} 3.8.0 para o processamento linguístico.

Uma adaptação análoga foi necessária para o CPC. O trabalho original ajusta um Mistral-7B-Instruct com o conjunto CQR para produzir representações de sentenças condicionadas ao contexto \cite{li2024cpc}. O repositório oficial disponibiliza o código, o conjunto CQR e compressores pré-treinados baseados em Mistral-7B e Llama-1B \cite{workday2025cpc}; ainda assim, acrescentar outro modelo desse porte à mesma máquina ampliaria a disputa de memória com as LLMs locais avaliadas. Na triagem experimental, o codificador é, portanto, substituído pelo all-MiniLM-L6-v2, da coleção Sentence Transformers no Hugging Face \cite{huggingface2022allminilm}. O compressor divide a partição em sentenças, calcula separadamente os \textit{embeddings} das sentenças e da pergunta, ordena as sentenças pela similaridade de cosseno e preserva as mais relevantes dentro da taxa de retenção, mantendo a ordem textual original. Essa variante é denominada CPC-MiniLM. Ela preserva a seleção extrativa em nível de sentença e o condicionamento à pergunta, mas não incorpora o treinamento contrastivo nem a representação contextual do CPC original. Consequentemente, trata-se de uma aproximação operacional, e diferenças de qualidade e latência entre CPC-MiniLM e os demais compressores não devem ser atribuídas diretamente ao método publicado por Liskavets et al. \cite{li2024cpc}.

\paragraph{Rastreabilidade da triagem.}
As quatro estratégias adicionadas são executadas somente quando nomeadas explicitamente e não foram incluídas no grupo \texttt{all}; essa decisão preserva a reprodutibilidade da etapa inicial da matriz local. A triagem usa o ambiente \texttt{.venv}, que contém as dependências dos quatro compressores, e grava seus artefatos em \path{tests/ollama-local/benchmark_matrix_compressores}, separado de \path{benchmark_matrix_artigo}. O protocolo completo prevê os cinco modelos, os orçamentos de 8192 e 32768 \textit{tokens}, os sete \textit{benchmarks} e a taxa de retenção de 0,5. O modo \texttt{quick} materializa três casos de resposta a perguntas e um caso de sumarização no LongBench, enquanto cada um dos outros seis arquivos locais contém três exemplos, totalizando 22 casos por estratégia e 88 resultados por combinação de modelo e orçamento.

Os artefatos consolidados reúnem as dez combinações de modelo e orçamento, correspondentes aos cinco modelos avaliados com 8192 e 32768 \textit{tokens}. Cada combinação registra 88 execuções, resultantes da aplicação dos quatro compressores aos mesmos 22 casos, totalizando 880 observações. Uma geração de resposta do CPC-MiniLM com o DeepSeek-R1 32B em 8192 \textit{tokens}, no caso \texttt{meeting\_summarization\_0}, terminou por \textit{timeout}. A saída de erro não foi tratada como resposta e foi excluída das médias de pontuação e latência, de modo que essas duas métricas reúnem 879 respostas válidas; a taxa retida e a fidelidade dessa execução permanecem válidas, pois a compressão foi concluída antes da chamada ao modelo alvo. A execução utiliza um processo por combinação de modelo e orçamento, no qual as quatro estratégias são registradas conjuntamente; assim, os auxiliares GPT-2 (LLMLingua-GPT2), XLM-RoBERTa (LLMLingua-2), GPT-2 (Selective Context) e MiniLM (CPC-MiniLM) podem coexistir durante a rodada. O processo é encerrado antes da combinação seguinte, liberando o modelo alvo e os auxiliares. Essa organização preserva a comparabilidade de qualidade dentro de cada combinação, mas as latências devem ser interpretadas como custo descritivo de ponta a ponta no ambiente local, e não como medição isolada do custo de cada compressor.

O orçamento é derivado independentemente para as janelas de 8192 e 32768 \textit{tokens}. Em cada rodada são descontados o \textit{template} completo, a pergunta, a reserva de saída e uma margem de segurança para diferenças entre o contador \textit{tiktoken} usado nas métricas e o tokenizador do modelo alvo. A implementação registra o custo previsto pelo resolvedor, o custo recontado do contexto final, o orçamento não utilizado e eventuais compressores indisponíveis, permitindo auditar cada execução.

\subsection{Aplicação da Solução de Ponta a Ponta}
\label{subsec:aplicacao}

Esta subseção rastreia um caso concreto por todos os estágios do \textit{framework}, de modo a tornar explícito o funcionamento descrito nas subseções anteriores. O caso provém do HotpotQA e foi processado sob a configuração principal, com o Llama 3.1 8B e orçamento de 8192 \textit{tokens}. A pergunta é ``Are the Laleli Mosque and Esma Sultan Mansion located in the same neighborhood?'', e o contexto reúne passagens sobre entidades otomanas, das quais apenas uma contém a informação necessária para a resposta.

O particionamento estrutural produz três partições disjuntas do tipo \textit{paragraph\_group}. A função de importância pontua cada partição pela combinação de relevância à pergunta, densidade de informação e saliência posicional, com os pesos 0,6, 0,2 e 0,2 definidos na Seção~\ref{subsec:importancia}. A Tabela~\ref{tab:exemplo-particoes} apresenta a decomposição da importância. A partição 2 reúne os verbetes da Laleli Mosque e da Esma Sultan Mansion e recebe relevância e densidade máximas, o que a torna a partição de maior importância. A partição 1 obtém saliência posicional máxima por ocupar a região central da entrada, enquanto a partição 0 apresenta relevância intermediária.

\begin{table}[htbp]
\centering
\caption{Particionamento e decomposição da importância no caso analisado}
\label{tab:exemplo-particoes}
\begin{tabular}{cccccc}
\toprule
\textbf{Partição} & \textbf{Caracteres} & \textbf{Relevância} & \textbf{Densidade} & \textbf{Saliência} & \textbf{Importância} \\
\midrule
0 & 1201 & 0,384 & 0,000 & 0,000 & 0,231 \\
1 & 1394 & 0,000 & 0,313 & 1,000 & 0,263 \\
2 & 1013 & 1,000 & 1,000 & 0,000 & \textbf{0,800} \\
\bottomrule
\end{tabular}
\end{table}

Para cada partição, o gerador materializa as opções de compressão e calcula, por opção, o custo em \textit{tokens}, a fidelidade e o valor $Q = I(t_j;q)\,f_{j,o}$. O conjunto reúne a cópia integral, a omissão e os compressores CPC-MiniLM e LLMLingua-2 nas taxas de 0,7, 0,5 e 0,3. A Tabela~\ref{tab:exemplo-opcoes} lista essas opções e destaca a escolhida pelo resolvedor em cada partição. Como a fidelidade multiplica a importância, a cópia integral concentra o maior valor nas três partições. As opções comprimidas ilustram, contudo, a heterogeneidade dos compressores descrita na Seção~\ref{subsec:pool}. Na partição 2, o LLMLingua-2 com taxa 0,7 alcança fidelidade de 0,843 e valor de 0,674, enquanto o CPC-MiniLM na mesma taxa obtém fidelidade de 0,500 e valor de 0,400. Sob pressão de orçamento, essa diferença de valor por \textit{token} leva o resolvedor a preferir o LLMLingua-2 nessa partição.

{
\footnotesize
\setlength{\tabcolsep}{6pt}
\begin{longtable}{lrrrr}
\caption{Opções de compressão avaliadas por partição, com a escolha do resolvedor em destaque}
\label{tab:exemplo-opcoes} \\
\toprule
\textbf{Compressor} & \textbf{Taxa} & \textbf{Custo} & \textbf{Fidelidade} & \textbf{Valor} \\
\midrule
\endfirsthead
\toprule
\textbf{Compressor} & \textbf{Taxa} & \textbf{Custo} & \textbf{Fidelidade} & \textbf{Valor} \\
\midrule
\endhead
\bottomrule
\endfoot
\bottomrule
\endlastfoot
\multicolumn{5}{l}{\textit{Partição 0 (importância 0,231)}} \\*
Cópia integral & 1,00 & 333 & 1,000 & \textbf{0,231} \\*
CPC-MiniLM     & 0,70 & 234 & 0,569 & 0,131 \\*
CPC-MiniLM     & 0,50 & 179 & 0,500 & 0,115 \\*
CPC-MiniLM     & 0,30 & 153 & 0,500 & 0,115 \\*
LLMLingua-2    & 0,70 & 233 & 0,610 & 0,141 \\*
LLMLingua-2    & 0,50 & 163 & 0,592 & 0,137 \\*
LLMLingua-2    & 0,30 & 104 & 0,479 & 0,111 \\*
Omissão        & 0,00 & 0   & 0,000 & 0,000 \\
\midrule
\multicolumn{5}{l}{\textit{Partição 1 (importância 0,263)}} \\*
Cópia integral & 1,00 & 369 & 1,000 & \textbf{0,263} \\*
CPC-MiniLM     & 0,70 & 249 & 0,308 & 0,081 \\*
CPC-MiniLM     & 0,50 & 176 & 0,308 & 0,081 \\*
CPC-MiniLM     & 0,30 & 121 & 0,308 & 0,081 \\*
LLMLingua-2    & 0,70 & 252 & 0,780 & 0,205 \\*
LLMLingua-2    & 0,50 & 179 & 0,698 & 0,183 \\*
LLMLingua-2    & 0,30 & 110 & 0,500 & 0,131 \\*
Omissão        & 0,00 & 0   & 0,000 & 0,000 \\
\midrule
\multicolumn{5}{l}{\textit{Partição 2 (importância 0,800)}} \\*
Cópia integral & 1,00 & 279 & 1,000 & \textbf{0,800} \\*
CPC-MiniLM     & 0,70 & 191 & 0,500 & 0,400 \\*
CPC-MiniLM     & 0,50 & 137 & 0,100 & 0,080 \\*
CPC-MiniLM     & 0,30 & 116 & 0,100 & 0,080 \\*
LLMLingua-2    & 0,70 & 188 & 0,843 & 0,674 \\*
LLMLingua-2    & 0,50 & 136 & 0,795 & 0,636 \\*
LLMLingua-2    & 0,30 & 86  & 0,700 & 0,560 \\*
Omissão        & 0,00 & 0   & 0,000 & 0,000 \\
\end{longtable}
}

A alocação resulta da resolução da mochila de múltipla escolha sobre essas opções. Neste caso, o custo da cópia integral das três partições soma 981 \textit{tokens}, valor inferior ao orçamento de 7390 \textit{tokens} disponível para o contexto. O resolvedor mantém, portanto, as três partições na cópia integral, e a penalização de transição não altera a escolha, pois todas as partições permanecem na mesma família. O contexto reconstruído concatena as três partições na ordem original e preserva a passagem que sustenta a resposta, na qual a Laleli Mosque é situada em Laleli, no distrito de Fatih, e a Esma Sultan Mansion no bairro de Ortaköy. O \textit{prompt} final combina esse contexto de 979 \textit{tokens} com a pergunta e a instrução de resposta.

O comportamento compressivo do \textit{framework} torna-se visível quando o orçamento é restritivo. Em um caso de contexto longo processado na mesma configuração, a cópia integral das 28 partições exigiria 9832 \textit{tokens}, acima do orçamento de 7396 \textit{tokens}. Nessa situação, o resolvedor manteve 16 partições na cópia integral, comprimiu 8 com o LLMLingua-2 e 4 com o CPC-MiniLM, e utilizou 7355 \textit{tokens}. A alocação preserva as partições de maior importância e transfere a redução para as demais, o que caracteriza a política seletiva por partição implementada pelo \textit{framework}.

\section{Perguntas de Pesquisa e Hipóteses}
\label{sec:perguntas-hipoteses}

A avaliação do \textit{framework} é orientada pelas seguintes perguntas de pesquisa.

\begin{itemize}[leftmargin=1.6em]
\item \textbf{PP1}. A alocação seletiva de compressão por partição preserva mais informação útil do que a compressão uniforme sob o mesmo orçamento de \textit{tokens}?

\item \textbf{PP2}. A função de importância, ao combinar relevância à tarefa, densidade de informação e saliência posicional, melhora a qualidade das respostas em comparação com uma alocação de orçamento independente da posição?

\item \textbf{PP3}. A penalização de transição entre partições adjacentes melhora a coerência do \textit{prompt} reconstruído sem prejudicar a acurácia final?

\item \textbf{PP4}. Como o \textit{framework} se comporta em diferentes escalas de modelo e orçamentos de contexto?
\end{itemize}

A partir dessas perguntas, formulam-se três hipóteses de trabalho.

\begin{itemize}[leftmargin=1.6em]
\item \textbf{H1}. Sob um orçamento fixo, a alocação seletiva por partição preserva mais informação relevante para a tarefa do que a compressão uniforme, resultando em acurácia igual ou superior nos \textit{benchmarks} de contexto longo.

\item \textbf{H2}. A inclusão da saliência posicional na função de importância reduz o efeito do viés \textit{lost in the middle}, ao proteger regiões centrais de alto valor da compressão excessiva.

\item \textbf{H3}. A penalização de transição reduz artefatos de incoerência entre trechos adjacentes sem perda significativa de acurácia.
\end{itemize}

As hipóteses definem critérios de avaliação e não são tratadas como resultados já obtidos.

\section{Plano Experimental para Avaliação do \textit{Framework}}
\label{sec:plano-experimental}

A avaliação do \textit{framework} é organizada em quatro componentes.

\begin{enumerate}[leftmargin=1.8em]
\item \textbf{\textit{Benchmarks}}. Utilização da mesma suíte de contexto longo empregada na validação empírica, descrita na Seção~\ref{subsec:desenho}, para manter a comparabilidade com a matriz de técnicas uniformes.

\item \textbf{\textit{Baselines}}. Comparação com a estratégia \textit{Raw}, com a \textit{Sliding Window} uniforme e com a aplicação uniforme de cada compressor candidato, isolando o efeito da alocação seletiva.

\item \textbf{Métricas}. Registro de acurácia, latência e taxa de compressão efetiva, complementadas por medidas de preservação de informação e fidelidade semântica.

\item \textbf{Estudos de ablação}. Remoção sucessiva de componentes, incluindo a saliência posicional da função de importância e a penalização de transição, para estimar a contribuição individual de cada componente.
\end{enumerate}

O plano prioriza a preservação de informação porque o objetivo do \textit{framework} é alocar o orçamento segundo a relevância das partições, e não apenas reduzir o número total de \textit{tokens}. Uma primeira execução desse plano é apresentada no Capítulo~\ref{cap:resultados}, restrita ao Llama 3.1 8B e a duas janelas de contexto, 8192 e 2048 \textit{tokens}, com as estratégias \textit{Raw}, MCKP e controle uniforme sobre os sete \textit{benchmarks}. A execução sobre os demais modelos e os estudos de ablação da função de importância e da penalização de transição permanecem como trabalho futuro no Capítulo~\ref{cap:conclusao}.
